\section{Supports and commutants of local observables}

\begin{lemma}[Disjoint finite-region observables commute]
    \label{thm:qca_disjoint_local_inclusion_commute}
    \lean{SpinChain.localInclusion_commute_of_disjoint}
    \leanok
    \uses{def:qca_local_inclusion}
    If $\Lambda\cap\Gamma=\varnothing$, $A\in\mathcal A_\Lambda$, and
    $B\in\mathcal A_\Gamma$, then their images in the algebra of the union
    satisfy
    \begin{align}
        \iota_{\Lambda,\Lambda\cup\Gamma}(A)
        \iota_{\Gamma,\Lambda\cup\Gamma}(B)
         & =\iota_{\Gamma,\Lambda\cup\Gamma}(B)
        \iota_{\Lambda,\Lambda\cup\Gamma}(A).
    \end{align}
    In \cite[lines~2292--2300]{Cirac2017MPU}, this identity provides the local commutativity of disjoint regions in the
    quasi-local algebra.
    % Source context: paper_v2.tex 2292--2300.
\end{lemma}

\begin{proof}\leanok
    For configurations $x,y$ on $\Lambda\cup\Gamma$, split each intermediate
    configuration $z$ into $(z|_\Lambda,z|_\Gamma)$.  The matrix entry of the
    left-hand side is
    \begin{align}
        \sum_{z_\Lambda,z_\Gamma}
        A(x|_\Lambda,z_\Lambda)
        \delta_{x|_\Gamma,z_\Gamma}
        B(z_\Gamma,y|_\Gamma)
        \delta_{z_\Lambda,y|_\Lambda}
         & =A(x|_\Lambda,y|_\Lambda)
        B(x|_\Gamma,y|_\Gamma).
    \end{align}
    Reversing the two factors gives the same scalar product.
\end{proof}

\begin{lemma}[Finite relative commutant]
    \label{thm:qca_finite_relative_commutant}
    \lean{SpinChain.mem_range_localInclusion_iff_commute_complement}
    \leanok
    \uses{def:qca_local_inclusion}
    Let $\Lambda\subseteq\Gamma$ be finite regions and
    $A\in\mathcal A_\Gamma$. Then $A$ acts trivially outside $\Lambda$ if
    and only if it commutes with the full observable algebra on the complement:
    \begin{align}
        A\in\iota_{\Lambda,\Gamma}(\mathcal A_\Lambda)
         & \iff
        \forall B\in\mathcal A_{\Gamma\setminus\Lambda},\quad
        A\iota_{\Gamma\setminus\Lambda,\Gamma}(B)
        =\iota_{\Gamma\setminus\Lambda,\Gamma}(B)A.
    \end{align}
    Unlike the results of Schumacher--Werner or the QCA appendix of \cite{Cirac2017MPU}, this lemma is a finite-dimensional tensor-factor
    identity.
\end{lemma}

\begin{proof}\leanok
    \uses{lem:amplified_center_scalar}
    Splitting configurations gives
    $\mathcal A_\Gamma\cong
        \mathcal A_\Lambda\otimes\mathcal A_{\Gamma\setminus\Lambda}$. Under this
    identification, the two inclusions are $C\mapsto C\otimes\mathbf 1$ and
    $B\mapsto\mathbf 1\otimes B$. For an arbitrary matrix written in blocks indexed
    by configurations on $\Lambda$, commuting with every
    $\mathbf 1\otimes B$ forces each block to commute with the full matrix algebra on
    $\Gamma\setminus\Lambda$ and is therefore scalar. Thus the matrix is
    $C\otimes\mathbf 1$ for some $C\in\mathcal A_\Lambda$. The converse follows
    from
    \begin{align}
        (C\otimes\mathbf 1)(\mathbf 1\otimes B)
         & =C\otimes B
        =(\mathbf 1\otimes B)(C\otimes\mathbf 1).
    \end{align}
\end{proof}

\subsection{Bipartite support algebras}

Let $I$ and $J$ be finite sets.  For
$X\in M_{I\times J}(\mathbb C)$, the existing left coefficient slice and
right operator block are the matrices
\begin{align}
    X^{\mathrm L}_{rs}(i,j) & =X((i,r),(j,s)), &
    X^{\mathrm R}_{ij}(r,s) & =X((i,r),(j,s)).
\end{align}
For a star-subalgebra $\mathcal D\subseteq M_K(\mathbb C)$, write
$\mathcal D'=\{Y\in M_K(\mathbb C)\mid [Y,Z]=0\text{ for every }Z\in\mathcal D\}$
for its commutant.  The construction below is the full complex matrix-algebra
specialization of the coefficient construction in
\cite[Section~4.3]{SchumacherWerner2004QCA} and in \cite[Lemma~7]{Gross2012IndexTheory} of the leastness and commutant
statements.

The commutant characterization of \cite[Lemma~7]{Gross2012IndexTheory} requires
star-closure \cite[lines~1161--1171]{SchumacherWerner2004QCA}, as documented
in \cite{gap:gnvw12_support_algebra_star_closure}. For the MPU spin chain, full
matrix factors $\mathcal B_1=M_I(\mathbb C)$ and $\mathcal B_2=M_J(\mathbb C)$
as in \cite[lines~2320--2334]{Cirac2017MPU} and \cite{gap:gnvw12_support_algebra_full_matrix_scope}
suffice.

\begin{lemma}[Adjoint of a right coefficient block]
    \label{thm:qca_operator_block_adjoint}
    \lean{Matrix.operatorBlock_conjTranspose}
    \leanok
    For every $X\in M_{I\times J}(\mathbb C)$ and $i,j\in I$,
    \begin{align}
        (X^\dagger)^{\mathrm R}_{ij}
         & =\bigl(X^{\mathrm R}_{ji}\bigr)^\dagger.
    \end{align}
\end{lemma}

\begin{proof}\leanok
    At the $(r,s)$ entry, both sides equal
    $\overline{X((j,s),(i,r))}$.
\end{proof}

\begin{definition}[Left coefficient set]
    \label{def:qca_left_coefficient_set}
    \lean{Matrix.leftCoefficientSet}
    \leanok
    Let $\mathcal A\subseteq M_{I\times J}(\mathbb C)$ be a star-subalgebra.
    Its set of left coefficients is
    \begin{align}
        \operatorname{Coeff}_{\mathrm L}(\mathcal A)
         & =\bigl\{X^{\mathrm L}_{rs}\ \bigm|\
        X\in\mathcal A,\ r,s\in J\bigr\}
        \subseteq M_I(\mathbb C).
    \end{align}
\end{definition}

\begin{definition}[Right coefficient set]
    \label{def:qca_right_coefficient_set}
    \lean{Matrix.rightCoefficientSet}
    \leanok
    Let $\mathcal A\subseteq M_{I\times J}(\mathbb C)$ be a star-subalgebra.
    Its set of right coefficients is
    \begin{align}
        \operatorname{Coeff}_{\mathrm R}(\mathcal A)
         & =\bigl\{X^{\mathrm R}_{ij}\ \bigm|\
        X\in\mathcal A,\ i,j\in I\bigr\}
        \subseteq M_J(\mathbb C).
    \end{align}
\end{definition}

\begin{definition}[Left support algebra]
    \label{def:qca_left_support_algebra}
    \lean{Matrix.leftSupportAlgebra}
    \leanok
    \uses{def:qca_left_coefficient_set}
    The left support algebra of $\mathcal A$ is the star-subalgebra
    \begin{align}
        \operatorname{Spp}_{\mathrm L}(\mathcal A)
         & =\operatorname{alg}^*
        \bigl(\operatorname{Coeff}_{\mathrm L}(\mathcal A)\bigr)
        \subseteq M_I(\mathbb C)
    \end{align}
    generated by all left coefficients.
\end{definition}

\begin{definition}[Right support algebra]
    \label{def:qca_right_support_algebra}
    \lean{Matrix.rightSupportAlgebra}
    \leanok
    \uses{def:qca_right_coefficient_set}
    The right support algebra of $\mathcal A$ is the star-subalgebra
    \begin{align}
        \operatorname{Spp}_{\mathrm R}(\mathcal A)
         & =\operatorname{alg}^*
        \bigl(\operatorname{Coeff}_{\mathrm R}(\mathcal A)\bigr)
        \subseteq M_J(\mathbb C)
    \end{align}
    generated by all right coefficients.
\end{definition}

\begin{lemma}[Left coefficient generators belong to the support algebra]
    \label{thm:qca_left_coefficient_mem_support}
    \lean{Matrix.bipartite_slice_mem_left_support_algebra}
    \leanok
    \uses{def:qca_left_support_algebra}
    If $X\in\mathcal A$ and $r,s\in J$, then
    \begin{align}
        X^{\mathrm L}_{rs}
         & \in\operatorname{Spp}_{\mathrm L}(\mathcal A).
    \end{align}
\end{lemma}

\begin{proof}\leanok
    \uses{def:qca_left_coefficient_set, def:qca_left_support_algebra}
    The matrix $X^{\mathrm L}_{rs}$ is one of the generators of the defining
    star-algebra closure.
\end{proof}

\begin{lemma}[Right coefficient generators belong to the support algebra]
    \label{thm:qca_right_coefficient_mem_support}
    \lean{Matrix.operator_block_mem_right_support_algebra}
    \leanok
    \uses{def:qca_right_support_algebra}
    If $X\in\mathcal A$ and $i,j\in I$, then
    \begin{align}
        X^{\mathrm R}_{ij}
         & \in\operatorname{Spp}_{\mathrm R}(\mathcal A).
    \end{align}
\end{lemma}

\begin{proof}\leanok
    \uses{def:qca_right_coefficient_set, def:qca_right_support_algebra}
    The matrix $X^{\mathrm R}_{ij}$ is one of the generators of the defining
    star-algebra closure.
\end{proof}

\begin{lemma}[Reconstruction from left coefficients]
    \label{thm:qca_left_coefficient_reconstruction}
    \lean{Matrix.eq_sum_left_coefficient_kronecker_single}
    \leanok
    For every $X\in M_{I\times J}(\mathbb C)$,
    \begin{align}
        X & =\sum_{r,s\in J}X^{\mathrm L}_{rs}\otimes E_{rs}.
    \end{align}
\end{lemma}

\begin{proof}\leanok
    At the entry $((i,r),(j,s))$, only the matrix unit $E_{rs}$ contributes,
    and its coefficient is $X^{\mathrm L}_{rs}(i,j)=X((i,r),(j,s))$.
\end{proof}

\begin{lemma}[Reconstruction from right coefficients]
    \label{thm:qca_right_coefficient_reconstruction}
    \lean{Matrix.eq_sum_single_kronecker_right_coefficient}
    \leanok
    For every $X\in M_{I\times J}(\mathbb C)$,
    \begin{align}
        X & =\sum_{i,j\in I}E_{ij}\otimes X^{\mathrm R}_{ij}.
    \end{align}
\end{lemma}

\begin{proof}\leanok
    At the entry $((i,r),(j,s))$, only the matrix unit $E_{ij}$ contributes,
    and its coefficient is $X^{\mathrm R}_{ij}(r,s)=X((i,r),(j,s))$.
\end{proof}

\begin{lemma}[Minimality of the left support algebra]
    \label{thm:qca_left_support_minimality}
    \lean{Matrix.left_support_algebra_le_iff}
    \leanok
    \uses{def:qca_left_support_algebra}
    For every star-subalgebra $\mathcal B\subseteq M_I(\mathbb C)$,
    \begin{align}
        \operatorname{Spp}_{\mathrm L}(\mathcal A)\subseteq\mathcal B
        \quad\Longleftrightarrow\quad
         & X^{\mathrm L}_{rs}\in\mathcal B
        \quad\text{for every }X\in\mathcal A\text{ and }r,s\in J.
    \end{align}
\end{lemma}

\begin{proof}\leanok
    \uses{def:qca_left_coefficient_set, def:qca_left_support_algebra}
    The universal property of the star-subalgebra generated by
    $\operatorname{Coeff}_{\mathrm L}(\mathcal A)$ gives the equivalence.
\end{proof}

\begin{lemma}[Minimality of the right support algebra]
    \label{thm:qca_right_support_minimality}
    \lean{Matrix.right_support_algebra_le_iff}
    \leanok
    \uses{def:qca_right_support_algebra}
    For every star-subalgebra $\mathcal C\subseteq M_J(\mathbb C)$,
    \begin{align}
        \operatorname{Spp}_{\mathrm R}(\mathcal A)\subseteq\mathcal C
        \quad\Longleftrightarrow\quad
         & X^{\mathrm R}_{ij}\in\mathcal C
        \quad\text{for every }X\in\mathcal A\text{ and }i,j\in I.
    \end{align}
\end{lemma}

\begin{proof}\leanok
    \uses{def:qca_right_coefficient_set, def:qca_right_support_algebra}
    The universal property of the star-subalgebra generated by
    $\operatorname{Coeff}_{\mathrm R}(\mathcal A)$ gives the equivalence.
\end{proof}

\begin{definition}[Kronecker-product submodule]
    \label{def:qca_kronecker_submodule}
    \lean{Matrix.kroneckerSubmodule}
    \leanok
    For submodules $S\subseteq M_I(\mathbb C)$ and
    $T\subseteq M_J(\mathbb C)$, define
    \begin{align}
        S\boxtimes T
         & =\operatorname{span}_{\mathbb C}
        \{A\otimes B\mid A\in S,\ B\in T\}
        \subseteq M_{I\times J}(\mathbb C).
    \end{align}
    This is the tensor-product subspace used below, rather than a
    conjunction of two one-sided containment conditions.
\end{definition}

\begin{lemma}[Coefficient criterion for the Kronecker-product submodule]
    \label{thm:qca_mem_kronecker_submodule_iff}
    \lean{Matrix.mem_kroneckerSubmodule_iff}
    \leanok
    \uses{def:qca_kronecker_submodule}
    For $X\in M_{I\times J}(\mathbb C)$,
    \begin{align}
        X\in S\boxtimes T
        \quad\Longleftrightarrow\quad
         & X^{\mathrm L}_{rs}\in S
        \quad\text{for every }r,s\in J, \\
         & \text{and}\quad
        X^{\mathrm R}_{ij}\in T
        \quad\text{for every }i,j\in I.
    \end{align}
\end{lemma}

\begin{proof}\leanok
    \uses{def:qca_kronecker_submodule}
    For $A\in S$ and $B\in T$, direct evaluation gives
    \begin{align}
        (A\otimes B)^{\mathrm L}_{rs} & =B_{rs}A, &
        (A\otimes B)^{\mathrm R}_{ij} & =A_{ij}B.
    \end{align}
    Thus the forward implication extends by linearity.
    Conversely, expand the right blocks in a basis of $T$.  Applying the dual
    coordinate functionals to those blocks produces matrices in $S$ from the
    left-slice hypothesis, and the resulting finite sum reconstructs $X$ as a
    sum of Kronecker products from $S$ and $T$.
\end{proof}

\begin{lemma}[Kronecker-submodule criterion for left support]
    \label{thm:qca_left_support_tensor_submodule}
    \lean{Matrix.left_support_algebra_le_iff_le_kronecker_submodule}
    \leanok
    \uses{def:qca_left_support_algebra, def:qca_kronecker_submodule}
    For every star-subalgebra $\mathcal B\subseteq M_I(\mathbb C)$,
    \begin{align}
        \operatorname{Spp}_{\mathrm L}(\mathcal A)\subseteq\mathcal B
        \quad\Longleftrightarrow\quad
        \mathcal A & \subseteq
        \mathcal B\boxtimes M_J(\mathbb C).
    \end{align}
\end{lemma}

\begin{proof}\leanok
    \uses{thm:qca_left_support_minimality, thm:qca_mem_kronecker_submodule_iff}
    By Lemma~\ref{thm:qca_mem_kronecker_submodule_iff} with the right
    submodule the full matrix space, membership in the Kronecker span is
    exactly membership of every left coefficient in $\mathcal B$, so
    Lemma~\ref{thm:qca_left_support_minimality} gives the equivalence.
\end{proof}

\begin{lemma}[Kronecker-submodule criterion for right support]
    \label{thm:qca_right_support_kronecker_submodule}
    \lean{Matrix.right_support_algebra_le_iff_le_kronecker_submodule}
    \leanok
    \uses{def:qca_right_support_algebra, def:qca_kronecker_submodule}
    For every star-subalgebra $\mathcal C\subseteq M_J(\mathbb C)$,
    \begin{align}
        \operatorname{Spp}_{\mathrm R}(\mathcal A)\subseteq\mathcal C
        \quad\Longleftrightarrow\quad
        \mathcal A & \subseteq M_I(\mathbb C)\boxtimes\mathcal C.
    \end{align}
\end{lemma}

\begin{proof}\leanok
    \uses{thm:qca_right_support_minimality, thm:qca_mem_kronecker_submodule_iff}
    The coefficient criterion reduces the right-hand containment to membership
    of every right block in $\mathcal C$.  This is equivalent to the left-hand
    inclusion by Lemma~\ref{thm:qca_right_support_minimality}.
\end{proof}

\begin{lemma}[Two-sided support-algebra containment]
    \label{thm:qca_two_sided_support_containment}
    \lean{Matrix.le_support_kroneckerSubmodule}
    \leanok
    \uses{def:qca_left_support_algebra, def:qca_right_support_algebra, def:qca_kronecker_submodule}
    Every bipartite star-subalgebra satisfies the two-sided containment
    \begin{align}
        \mathcal A
         & \subseteq
        \operatorname{Spp}_{\mathrm L}(\mathcal A)
        \boxtimes
        \operatorname{Spp}_{\mathrm R}(\mathcal A).
    \end{align}
    This is the full complex matrix-algebra specialization of the
    support-algebra containment in Schumacher--Werner Section~4.3 and of
    Gross--Nesme--Vogts--Werner equation~(22).  The latter follows after
    applying Lemma~7 to both tensor factors.  The displayed assertion is an
    inclusion in the tensor-product span itself, not merely the conjunction of
    the two one-sided inclusions.  In \cite[lines~2320--2334]{Cirac2017MPU}, the MPU appendix applies this construction
    to
    $\mathcal R_{2x}$.
\end{lemma}

\begin{proof}\leanok
    \uses{thm:qca_mem_kronecker_submodule_iff, thm:qca_left_coefficient_mem_support, thm:qca_right_coefficient_mem_support}
    If $X\in\mathcal A$, every left coefficient of $X$ lies in the left support
    algebra and every right coefficient lies in the right support algebra.
    Lemma~\ref{thm:qca_mem_kronecker_submodule_iff} places $X$ in their
    Kronecker-product submodule.
\end{proof}

\begin{lemma}[Left Kronecker commutation is coefficientwise]
    \label{thm:qca_left_kronecker_commute_iff}
    \lean{Matrix.commute_left_kronecker_embed_iff}
    \leanok
    For $B\in M_I(\mathbb C)$ and $X\in M_{I\times J}(\mathbb C)$,
    \begin{align}
        [B\otimes\mathbf 1_J,X]=0
        \quad\Longleftrightarrow\quad
        [B,X^{\mathrm L}_{rs}] & =0
        \quad\text{for every }r,s\in J.
    \end{align}
\end{lemma}

\begin{proof}\leanok
    For $i,j\in I$ and $r,s\in J$, the two product entries are
    \begin{align}
        [(B\otimes\mathbf 1_J)X]_{(i,r),(j,s)}
         & =\sum_{k\in I}B_{ik}X_{(k,r),(j,s)}
        =(BX^{\mathrm L}_{rs})_{ij},           \\
        [X(B\otimes\mathbf 1_J)]_{(i,r),(j,s)}
         & =\sum_{k\in I}X_{(i,r),(k,s)}B_{kj}
        =(X^{\mathrm L}_{rs}B)_{ij}.
    \end{align}
    Equality of the two products is therefore equivalent to
    $[B,X^{\mathrm L}_{rs}]=0$ for every $r,s\in J$.
\end{proof}

\begin{lemma}[Right Kronecker commutation is coefficientwise]
    \label{thm:qca_right_kronecker_commute_iff}
    \lean{Matrix.commute_right_kronecker_embed_iff}
    \leanok
    For $C\in M_J(\mathbb C)$ and $X\in M_{I\times J}(\mathbb C)$,
    \begin{align}
        [\mathbf 1_I\otimes C,X]=0
        \quad\Longleftrightarrow\quad
        [C,X^{\mathrm R}_{ij}] & =0
        \quad\text{for every }i,j\in I.
    \end{align}
\end{lemma}

\begin{proof}\leanok
    For $i,j\in I$ and $r,s\in J$, the two product entries are
    \begin{align}
        [(\mathbf 1_I\otimes C)X]_{(i,r),(j,s)}
         & =\sum_{t\in J}C_{rt}X_{(i,t),(j,s)}
        =(CX^{\mathrm R}_{ij})_{rs},           \\
        [X(\mathbf 1_I\otimes C)]_{(i,r),(j,s)}
         & =\sum_{t\in J}X_{(i,r),(j,t)}C_{ts}
        =(X^{\mathrm R}_{ij}C)_{rs}.
    \end{align}
    Equality of the two products is therefore equivalent to
    $[C,X^{\mathrm R}_{ij}]=0$ for every $i,j\in I$.
\end{proof}

\begin{lemma}[Left support-algebra centralizer]
    \label{thm:qca_left_support_centralizer}
    \lean{Matrix.left_kronecker_embed_mem_centralizer_iff}
    \leanok
    \uses{def:qca_left_support_algebra}
    For $B\in M_I(\mathbb C)$,
    \begin{align}
        B\otimes\mathbf 1_J\in\mathcal A'
        \quad\Longleftrightarrow\quad
        B & \in\operatorname{Spp}_{\mathrm L}(\mathcal A)'.
    \end{align}
    Equivalently,
    \begin{align}
        \operatorname{Spp}_{\mathrm L}(\mathcal A)'
         & =\{B\in M_I(\mathbb C)\mid B\otimes\mathbf 1_J\in\mathcal A'\}.
    \end{align}
\end{lemma}

\begin{proof}\leanok
    \uses{thm:qca_left_coefficient_mem_support, thm:qca_left_kronecker_commute_iff, def:qca_left_coefficient_set, def:qca_left_support_algebra}
    By Lemma~\ref{thm:qca_left_kronecker_commute_iff}, commutation of
    $B\otimes\mathbf 1_J$ with every $X\in\mathcal A$ is equivalent to
    commutation of $B$ with every left coefficient.  The coefficient set is
    closed under adjoints because $\mathcal A$ is a star-subalgebra.  The
    universal property of the generated star-subalgebra then extends this
    commutation relation to all of
    $\operatorname{Spp}_{\mathrm L}(\mathcal A)$; the converse follows from
    generator membership.
\end{proof}

\begin{lemma}[Right support-algebra centralizer]
    \label{thm:qca_right_support_centralizer}
    \lean{Matrix.right_kronecker_embed_mem_centralizer_iff}
    \leanok
    \uses{def:qca_right_support_algebra}
    For $C\in M_J(\mathbb C)$,
    \begin{align}
        \mathbf 1_I\otimes C\in\mathcal A'
        \quad\Longleftrightarrow\quad
        C & \in\operatorname{Spp}_{\mathrm R}(\mathcal A)'.
    \end{align}
    Equivalently,
    \begin{align}
        \operatorname{Spp}_{\mathrm R}(\mathcal A)'
         & =\{C\in M_J(\mathbb C)\mid \mathbf 1_I\otimes C\in\mathcal A'\}.
    \end{align}
\end{lemma}

\begin{proof}\leanok
    \uses{thm:qca_operator_block_adjoint, thm:qca_right_coefficient_mem_support, thm:qca_right_kronecker_commute_iff, def:qca_right_coefficient_set, def:qca_right_support_algebra}
    By Lemma~\ref{thm:qca_right_kronecker_commute_iff}, commutation of
    $\mathbf 1_I\otimes C$ with every $X\in\mathcal A$ is equivalent to
    commutation of $C$ with every right coefficient.  Lemma~\ref{thm:qca_operator_block_adjoint}
    shows that the right coefficient set is closed under adjoints.  The
    universal property of the generated star-subalgebra extends commutation to
    all of $\operatorname{Spp}_{\mathrm R}(\mathcal A)$; the converse follows
    from generator membership.
\end{proof}

\begin{lemma}[Commutation of overlapping support algebras]
    \label{thm:qca_overlapping_support_algebras_commute}
    \lean{Matrix.supportAlgebras_commute_of_overlappingLifts_commute}
    \leanok
    \uses{def:qca_left_support_algebra, def:qca_right_support_algebra}
    Let $I$, $J$, and $K$ be finite sets, and let
    \begin{align}
        \mathcal A_1 & \subseteq M_{I\times J}(\mathbb C), &
        \mathcal A_2 & \subseteq M_{J\times K}(\mathbb C)
    \end{align}
    be star-subalgebras.  Suppose that, for every
    $X\in\mathcal A_1$ and $Y\in\mathcal A_2$,
    \begin{align}
        (X\otimes\mathbf 1_K)(\mathbf 1_I\otimes Y)
         & =(\mathbf 1_I\otimes Y)(X\otimes\mathbf 1_K).
        \label{eq:qca_overlapping_support_lifts}
    \end{align}
    Then the two support algebras on the common factor commute:
    \begin{align}
        RL & =LR
           &     & \text{for every }
        R\in\operatorname{Spp}_{\mathrm R}(\mathcal A_1),\quad
        L\in\operatorname{Spp}_{\mathrm L}(\mathcal A_2).
        \label{eq:qca_overlapping_support_conclusion}
    \end{align}
    This is the full complex matrix-algebra specialization of
    Schumacher--Werner Lemma \texttt{sppcomm}
    \cite[lines~1174--1194]{SchumacherWerner2004QCA} and of \cite[Section~7.1, lines~1221--1246]{Gross2012IndexTheory} (GNVW Lemma
    \texttt{sppcomm})
    for commuting support algebras.

    The input star-subalgebra must be closed under involution, and the basis
    element belongs to the third factor; see the notes
    \cite{gap:gnvw12_support_algebra_star_closure} and
    \cite{gap:gnvw12_support_algebra_outer_factor_basis}. The statement treats
    full complex matrix algebras, which as documented in \cite[lines~2300--2306, 2313--2334]{Cirac2017MPU} and \cite{gap:gnvw12_support_algebra_full_matrix_scope} suffices for the MPU
    application to homogeneous spin
    chains.
\end{lemma}

\begin{proof}\leanok
    \uses{def:qca_left_coefficient_set, def:qca_right_coefficient_set, def:qca_left_support_algebra, def:qca_right_support_algebra}
    For $X\in\mathcal A_1$, $Y\in\mathcal A_2$, $i,i'\in I$, and
    $k,k'\in K$, applying
    Theorem~\cite[``Middle-factor coefficients of commuting overlapping operators'']{QICLean2026Blueprint} to
    (\ref{eq:qca_overlapping_support_lifts}) gives
    \begin{align}
        X^{\mathrm R}_{ii'}Y^{\mathrm L}_{kk'}
         & =Y^{\mathrm L}_{kk'}X^{\mathrm R}_{ii'}.
        \label{eq:qca_overlapping_support_coefficients}
    \end{align}
    Applying the same identity to $Y^\dagger\in\mathcal A_2$, with $k$ and
    $k'$ exchanged, also gives commutation with
    $(Y^{\mathrm L}_{kk'})^\dagger$.  Hence every generator of
    $\operatorname{Spp}_{\mathrm R}(\mathcal A_1)$ lies in the
    star-subalgebra commuting with all left coefficients of $\mathcal A_2$.
    Passing to the generated star-subalgebra yields
    \begin{align}
        LR          & =RL,
        \label{eq:qca_overlapping_support_first_closure} \\
        L^\dagger R & =RL^\dagger
        \notag
    \end{align}
    for every $R\in\operatorname{Spp}_{\mathrm R}(\mathcal A_1)$ and every
    left coefficient $L$ of $\mathcal A_2$.

    Since $\operatorname{Spp}_{\mathrm R}(\mathcal A_1)$ is closed under
    adjoints, applying~(\ref{eq:qca_overlapping_support_first_closure}) to $R^\dagger$ shows that
    each left coefficient commutes with both $R$ and $R^\dagger$.  Passing to
    the star-subalgebra generated by the left coefficients gives
    (\ref{eq:qca_overlapping_support_conclusion}).
\end{proof}

\subsection{Products of commuting matrix factors}

Let $\mathcal C$ be a nonzero complex star-algebra and let
$\mathcal A,\mathcal B\subseteq\mathcal C$ be unital star-subalgebras such
that $ab=ba$ for every $a\in\mathcal A$ and $b\in\mathcal B$.

\begin{definition}[Multiplication of commuting factors]
    \label{def:qca_commuting_factor_multiplication}
    \lean{StarSubalgebra.commutingMulMap,
        StarSubalgebra.commutingMulMap_tmul}
    \leanok
    Multiplication defines a unital star-algebra homomorphism
    \begin{align}
        \mu_{\mathcal A,\mathcal B}
        \colon \mathcal A\otimes_{\mathbb C}\mathcal B
         & \longrightarrow \mathcal C, \notag \\
        \mu_{\mathcal A,\mathcal B}(a\otimes b)
         & =ab.
    \end{align}
\end{definition}

\begin{lemma}[Algebraic product of commuting factors]
    \label{thm:qca_commuting_factor_product}
    \lean{StarSubalgebra.commutingMulMap_range,
        StarSubalgebra.sup_toSubmodule_eq_mul}
    \leanok
    \uses{def:qca_commuting_factor_multiplication}
    The image of $\mu_{\mathcal A,\mathcal B}$ is the algebraic join
    $\mathcal A\vee\mathcal B$.  Its underlying complex vector space is
    \begin{align}
        \mathcal A\mathcal B
         & =\operatorname{span}_{\mathbb C}
        \{ab\mid a\in\mathcal A,\ b\in\mathcal B\}.
        \label{eq:qca_commuting_factor_product_span}
    \end{align}
    Thus every element of the join is a finite complex-linear sum of
    products $ab$.
\end{lemma}

\begin{proof}\leanok
    \uses{def:qca_commuting_factor_multiplication}
    The involution reverses products.  Since $a^*\in\mathcal A$ and
    $b^*\in\mathcal B$ commute,
    \begin{align}
        \mu_{\mathcal A,\mathcal B}\bigl((a\otimes b)^*\bigr)
         & =a^*b^*=b^*a^*=(ab)^*.
    \end{align}
    The image contains both factors because $a=a\mathbf 1$ and
    $b=\mathbf 1b$.  Conversely, every pure tensor maps to an element of
    the join, and linearity gives the reverse inclusion.  The description
    of the underlying vector space is the range formula for the linear
    multiplication map.
\end{proof}

\begin{lemma}[Product of two full matrix factors]
    \label{thm:qca_commuting_full_matrix_product}
    \lean{StarSubalgebra.commutingMulMap_injective_of_equiv_matrix,
        StarSubalgebra.commutingSupEquivMatrix,
        StarSubalgebra.nonempty_equiv_sup_matrix,
        StarSubalgebra.finrank_sup_of_equiv_matrix}
    \leanok
    \uses{def:qca_commuting_factor_multiplication}
    Suppose that $\mathcal C=M_I(\mathbb C)$ for a finite nonempty set
    $I$, that $p,q>0$, and that
    \begin{align}
        \mathcal A & \cong M_p(\mathbb C), \notag \\
        \mathcal B & \cong M_q(\mathbb C)
    \end{align}
    as star-algebras.  Then $\mu_{\mathcal A,\mathcal B}$ is injective and
    \begin{align}
        \mathcal A\vee\mathcal B
         & \cong M_{pq}(\mathbb C), \notag \\
        \dim_{\mathbb C}(\mathcal A\vee\mathcal B)
         & =(pq)^2.
        \label{eq:qca_commuting_full_matrix_product}
    \end{align}
    This is the unital single-block specialization of Schumacher--Werner
    Proposition~\texttt{Cscom}
    \cite[lines~2119--2138]{SchumacherWerner2004QCA}, which in \cite[lines~1270--1308]{Gross2012IndexTheory} and \cite{gap:gnvw12_cscom_single_block_scope} is used for the adjacent
    support
    algebras.
\end{lemma}

\begin{proof}\leanok
    \uses{def:qca_commuting_factor_multiplication,
        thm:qca_commuting_factor_product}
    The source of the multiplication homomorphism is
    \begin{align}
        M_p(\mathbb C)\otimes_{\mathbb C}M_q(\mathbb C)
         & \cong M_{pq}(\mathbb C).
    \end{align}
    This is a simple ring.  The kernel of a unital homomorphism from a
    simple ring to a nonzero ring is not the whole ring, and is therefore
    zero.  Hence multiplication is injective.  Its image is the join by
    Lemma~\ref{thm:qca_commuting_factor_product}; the matrix presentation
    and dimension formula in
    Equation~(\ref{eq:qca_commuting_full_matrix_product}) follow.
\end{proof}

\begin{lemma}[Canonical Kronecker coordinates]
    \label{thm:qca_commuting_factor_kronecker_coordinates}
    \lean{Matrix.leftKroneckerEmbed_sup_rightKroneckerEmbed_toSubmodule}
    \leanok
    \uses{def:qca_kronecker_submodule}
    Let $I,J$ be finite sets, and let
    $\mathcal S\subseteq M_I(\mathbb C)$ and
    $\mathcal T\subseteq M_J(\mathbb C)$ be unital star-subalgebras.  In
    the canonical bipartite coordinates, the underlying complex vector
    space of the join of $\mathcal S\otimes\mathbf 1_J$ and
    $\mathbf 1_I\otimes\mathcal T$ is exactly
    \begin{align}
        \mathcal S\boxtimes\mathcal T
         & =\operatorname{span}_{\mathbb C}
        \{S\otimes T\mid S\in\mathcal S,\ T\in\mathcal T\}.
        \label{eq:qca_commuting_factor_kronecker_coordinates}
    \end{align}
\end{lemma}

\begin{proof}\leanok
    \uses{thm:qca_commuting_factor_product}
    If $I$ or $J$ is empty, the ambient matrix space is zero, so both
    subspaces in the claimed equality are zero.  Assume that both sets are
    nonempty.  The two canonical factors commute, and their products satisfy
    \begin{align}
        (S\otimes\mathbf 1_J)(\mathbf 1_I\otimes T)
         & =S\otimes T.
    \end{align}
    Equation~(\ref{eq:qca_commuting_factor_product_span}) therefore becomes
    exactly
    (\ref{eq:qca_commuting_factor_kronecker_coordinates}).
\end{proof}

\section{Periodic matrix product unitaries on a finite chain}

\begin{definition}[Finite-chain realization of a periodic MPU]
    \label{def:mpu_finite_chain_realization}
    \lean{SpinChain.finiteChainRegion,SpinChain.card_finiteChainRegion,
        SpinChain.finiteChainSiteEquiv,SpinChain.finiteChainConfigEquiv,
        SpinChain.finiteChainMPOMatrix}
    \leanok
    \uses{def:mpo_family, def:qca_local_algebra}
    For $a\in\mathbb Z$ and $N\in\mathbb N$, let
    $I_{a,N}=[a,a+N)\cap\mathbb Z$.  The increasing bijection
    $\{0,\ldots,N-1\}\simeq I_{a,N}$ identifies the periodic configuration
    basis with $\operatorname{Config}_d(I_{a,N})$.  In this basis, write
    $U_{a}^{(N)}$ for the matrix obtained by simultaneously reindexing the rows
    and columns of $U^{(N)}$.
\end{definition}

\begin{definition}[Residual words on an anchored interval]
    \label{def:mpu_anchored_residual_configurations}
    \lean{SpinChain.nextBlockPrefix,
        SpinChain.anchoredResidualConfig,
        SpinChain.anchoredResidualConfigEquiv,
        SpinChain.anchoredResidualConfigEquiv_apply}
    \leanok
    \uses{def:mpu_finite_chain_realization, def:mpdo_physical_blocking}
    Fix a positive block length $L$ and an integer anchor $a$. The number
    of sites from $a$ to the next globally aligned block boundary is
    $q(a,L)=(-a)\bmod L$. If an interval of length $N$ has the decomposition
    $N=p+mL+s$, its configuration may be specified by a prefix word of
    length $p$, a word of $m$ blocked letters, and a suffix word of length
    $s$. Their ordered concatenation defines a configuration on
    $I_{a,N}$, including when $m=0$. This concatenation is a bijection
    from triples of such words to configurations on the interval.
\end{definition}

\begin{lemma}[Decomposition of an arbitrary anchored interval]
    \label{thm:mpu_anchored_residual_decomposition}
    \lean{SpinChain.nextBlockPrefix_lt,
        SpinChain.nextBlockPrefix_aligned,
        SpinChain.exists_anchored_interval_decomposition,
        SpinChain.finiteChainConfigEquiv_apply_site,
        SpinChain.finiteChainRegion_sort,
        SpinChain.finiteChainSiteEquiv_val,
        SpinChain.anchoredResidualConfig_apply_site,
        SpinChain.middle_index_lt,
        SpinChain.aligned_middle_site_coord}
    \leanok
    \uses{def:mpu_anchored_residual_configurations,
        def:mpu_finite_chain_realization}
    For every $a\in\mathbb Z$, $N\in\mathbb N$, and $L>0$, there are
    $p,m,s\in\mathbb N$ such that
    \begin{align}
        p,s & <L, & N & =p+mL+s.
    \end{align}
    If $m>0$, then $a+p$ is divisible by $L$; the site at residue $r$ of
    the $t$th complete middle block has integer coordinate
    $a+p+tL+r$. If $N<q(a,L)$, then $p=N$ and $m=s=0$. The anchored
    configuration reads the prefix, middle blocks, and suffix in exactly
    this increasing site order.
\end{lemma}

\begin{proof}\leanok
    \uses{def:mpu_anchored_residual_configurations,
        def:mpu_finite_chain_realization}
    With $p=\min(N,q(a,L))$, dividing $N-p$ by $L$ with quotient
    $m$ and remainder $s$ gives $s<L$ and the
    length identity. If $m>0$, then $p=q(a,L)$ and $a+p$ is aligned.
    The order-preserving site identification sends position $k$ to $a+k$;
    applying it at $k=p+tL+r$ gives the middle-block coordinate. When
    $N<q(a,L)$, one has $p=N$ and both other lengths vanish.
\end{proof}

\begin{lemma}[Finite-chain MPO entries with residual sites]
    \label{thm:mpu_anchored_residual_mpo_entry}
    \lean{SpinChain.finiteChainMPOMatrix_apply_config,
        SpinChain.finiteChainMPOMatrix_apply_anchoredResidualConfig,
        SpinChain.finiteChainMPOMatrix_anchoredResidualTrace,
        SpinChain.anchoredResidualWordTrace,
        SpinChain.finiteChainMPOMatrix_anchoredResidualConfigEquiv}
    \leanok
    \uses{def:mpu_anchored_residual_configurations,
        def:mpu_finite_chain_realization, def:mpo_eval_word}
    Let $U$ be any MPO tensor and let $N=p+mL+s$. For two anchored
    configurations with residual prefix words $i^-,j^-$, blocked middle
    words $\sigma,\tau$, and residual suffix words $i^+,j^+$, the entry
    of the finite-chain matrix is
    \begin{align}
        (U_a^{(N)})_{i^-\widehat\sigma i^+,\,
                j^-\widehat\tau j^+}
         & =\operatorname{tr}\!\left(
        U^{i^-,j^-}(U^{[L]})^{\sigma,\tau}U^{i^+,j^+}\right).
    \end{align}
    This holds at every integer anchor and for arbitrary $L,p,m,s$,
    including an empty middle word. Neither unitarity nor divisibility
    of $N$ by $L$ is assumed.
\end{lemma}

\begin{proof}\leanok
    \uses{lem:mpdo_physical_blocking_residuals,
        def:mpu_finite_chain_realization}
    Undo the increasing reindexing of the finite-chain MPO matrix. The
    resulting periodic MPO entry is the trace of the prefix, blocked
    middle, and suffix products by residual-site blocking.
\end{proof}

\begin{definition}[Finite-chain MPU conjugation]
    \label{def:mpu_finite_chain_conjugation}
    \lean{SpinChain.finiteChainMPUUnitary,
        SpinChain.finiteChainConjugation,
        SpinChain.finiteChainConjugation_apply}
    \leanok
    \uses{def:mpu, def:mpu_finite_chain_realization}
    If $\mathcal U$ is an MPU and $N>1$, then $U_a^{(N)}$ is unitary and
    defines the star-algebra automorphism
    \begin{align}
        \operatorname{Ad}_{U_a^{(N)}}(B)
         & =U_a^{(N)}B\,U_a^{(N)\dagger}.
    \end{align}
    In \cite[lines~2300--2306]{Cirac2017MPU}, this conjugation uses the same factor
    order.
\end{definition}

\begin{proof}\leanok
    \uses{def:mpu, thm:mpu_reindex_unitary}
    Simultaneous reindexing preserves both unitarity equations for $U^{(N)}$.
    Conjugation by the resulting unitary preserves multiplication, the unit,
    complex scalars, and the adjoint.
\end{proof}

\begin{lemma}[Conjugation coefficients in residual-word coordinates]
    \label{thm:mpu_anchored_residual_conjugation_entry}
    \lean{SpinChain.finiteChainConjugation_localInclusion_apply_entry,
        SpinChain.finiteChainConjugation_localInclusion_anchoredResidualTrace}
    \leanok
    \uses{def:mpu_anchored_residual_configurations,
        thm:mpu_anchored_residual_mpo_entry,
        def:mpu_finite_chain_conjugation, def:qca_local_inclusion}
    Let $U$ be an MPU, $N=p+mL+s>1$, and $A\in\mathcal A_\Lambda$
    for $\Lambda\subseteq I_{a,N}$. Write $\widehat x$ for the anchored
    configuration represented by a triple $x=(x^-,x^0,x^+)$ of residual
    prefix, complete-block, and residual suffix words, and put
    \begin{align}
        T(x,y) & =\operatorname{tr}\!\left(
        U^{x^-,y^-}(U^{[L]})^{x^0,y^0}U^{x^+,y^+}\right).
    \end{align}
    For all such triples $x,y$, the matrix coefficient of finite-chain
    conjugation is
    \begin{align}
        \bigl(U_a^{(N)}\iota_{\Lambda,I_{a,N}}(A)
        U_a^{(N)\dagger}\bigr)_{\widehat x,\widehat y}
         & =\sum_{z,w} T(x,z)A_{\widehat z|_\Lambda,\widehat w|_\Lambda}
        \notag                                                           \\
         & \qquad\cdot
        \delta_{\widehat z|_{I_{a,N}\setminus\Lambda},
            \widehat w|_{I_{a,N}\setminus\Lambda}}
        \overline{T(y,w)}.
    \end{align}
    Both MPO factors retain their residual prefix and suffix products.
    This finite-chain identity does not assert independence of the
    interval or support of the conjugated observable.
\end{lemma}

\begin{proof}\leanok
    \uses{thm:mpu_anchored_residual_mpo_entry,
        def:mpu_finite_chain_conjugation, def:qca_local_inclusion}
    Expanding conjugation as a double sum over intermediate configurations shows that
    the local inclusion contributes the matrix entry of $A$ on $\Lambda$
    and the equality condition on its complement. Reindexing both sums by
    the residual-word bijection and applying the finite-chain residual-trace
    identity to each MPO coefficient completes the proof.
\end{proof}

\begin{definition}[Three-block coordinates on an anchored finite chain]
    \label{def:mpu_anchored_three_block_coordinates}
    \lean{MPOTensor.threeBlockOutputEquiv,
        MPOTensor.threeBlockOutputEquiv_apply,
        MPOTensor.threeBlockInputEquiv,
        MPOTensor.threeBlockInputEquiv_apply,
        MPOTensor.threeBlockOutputConfigEquiv,
        MPOTensor.threeBlockOutputConfigEquiv_apply,
        MPOTensor.threeBlockInputConfigEquiv,
        MPOTensor.threeBlockInputConfigEquiv_apply,
        MPOTensor.threeBlockFiniteChainMatrix}
    \leanok
    \uses{def:mpu_anchored_residual_configurations,
        def:mpu_finite_chain_realization,
        def:mpu_two_site_open_three_block}
    Fix $k,p,s\geq0$ and an integer anchor $a$. On the interval of length
    $N=p+3(2k)+s$, distinguish three consecutive complete blocks of
    length $2k$ between an arbitrary prefix and suffix. Let $H=d^k$.
    On the output side, the four neighboring half-blocks form a factor
    $X$ of dimension $H^4$; the prefix, suffix, and two exposed outer
    half-blocks form the exterior factor $E_{\mathrm{out}}$. On the input
    side, the central complete block forms a factor $J$ of dimension
    $H^2$; the prefix, suffix, and the two outer complete blocks form
    $E_{\mathrm{in}}$. Ordered concatenation gives bijections from
    $E_{\mathrm{out}}\times X$ and $E_{\mathrm{in}}\times J$ to the
    respective configuration bases of the same finite interval. Write
    $M_{a;p,k,s}$ for the finite-chain MPO matrix in these two bases.
\end{definition}

\begin{lemma}[Anchored three-block operator identity]
    \label{thm:mpu_anchored_three_block_operator}
    \lean{MPOTensor.threeBlockFiniteChainMatrix_slice,
        Matrix.intertwiner_of_slices,
        MPOTensor.threeBlockFiniteChainMatrix_intertwiner,
        MPOTensor.threeBlockFiniteChainMatrix_conjugation}
    \leanok
    \uses{def:mpu_anchored_three_block_coordinates,
        thm:mpu_original_site_three_block_window_entry,
        lem:mpu_two_site_open_three_block_locality,
        thm:mpu_anchored_residual_mpo_entry,
        def:mpu_finite_chain_conjugation}
    Suppose the two-site blocking of $U$ at length $k$ is equipped with
    standard-form factors $u,v$. For every observable $A$ on the central
    input block, let $B(A)$ be the corresponding observable on the two
    neighboring output links. Then
    \begin{align}
        M_{a;p,k,s}(I_{E_{\mathrm{in}}}\otimes A)
         & =(I_{E_{\mathrm{out}}}\otimes B(A))M_{a;p,k,s}.
        \label{eq:mpu_finite_chain_intertwiner}
    \end{align}
    This identity holds for arbitrary original-site prefix and suffix
    words; $U$ need not be an MPU. If $U$ is an MPU and $N>1$, then
    \begin{align}
        M_{a;p,k,s}(I_{E_{\mathrm{in}}}\otimes A)M_{a;p,k,s}^{\dagger}
         & =I_{E_{\mathrm{out}}}\otimes B(A).
        \label{eq:mpu_finite_chain_conjugation}
    \end{align}
    Thus the conjugated operator is the identity on every exterior
    output coordinate. This statement assumes the displayed given
    standard form and does not yet identify the right-hand side with
    a site-indexed local inclusion or prove independence of the
    enclosing interval.
\end{lemma}

\begin{proof}\leanok
    \uses{thm:mpu_original_site_three_block_window_entry,
        lem:mpu_two_site_open_three_block_locality,
        def:mpu_anchored_three_block_coordinates,
        def:mpu_finite_chain_conjugation}
    Fixing all exterior input and output coordinates makes a slice of
    $M_{a;p,k,s}$ equal to the original-site three-block window.
    The open-word identity intertwines each slice with the same $B(A)$.
    Summing over exterior coordinates yields~(\ref{eq:mpu_finite_chain_intertwiner}).
    When $U$ is an MPU, $M_{a;p,k,s}$ is unitary; postmultiplying~(\ref{eq:mpu_finite_chain_intertwiner}) by $M_{a;p,k,s}^\dagger$ yields~(\ref{eq:mpu_finite_chain_conjugation}).
\end{proof}

\begin{definition}[Finite-chain term on a local observable]
    \label{def:mpu_finite_chain_local_observable}
    \lean{SpinChain.finiteChainLocalObservable,
        SpinChain.finiteChainLocalObservable_apply}
    \leanok
    \uses{def:qca_local_inclusion, def:qca_algebraic_local_algebra, def:mpu_finite_chain_conjugation}
    Let $\Lambda\subseteq I_{a,N}$ and $A\in\mathcal A_\Lambda$.  Set
    $A_L=\iota_{\Lambda,I_{a,N}}(A)$.  For $N>1$, the finite-chain term in
    \cite[lines~2300--2306]{Cirac2017MPU} is the star-algebra map
    \begin{align}
        \mathcal A_\Lambda & \longrightarrow\mathcal A_{\mathrm{loc}}, \\
        A                  & \longmapsto
        \iota_{I_{a,N}}\!\left(U_a^{(N)}A_LU_a^{(N)\dagger}\right).
    \end{align}
    No stabilization in $N$ is asserted.
\end{definition}

\section{Locality and translation covariance}

\begin{lemma}[Disjointly supported local observables commute]
    \label{thm:qca_disjoint_supported_local_commute}
    \lean{SpinChain.localObservable_commute_of_disjoint,
        SpinChain.SupportedIn.commute_of_disjoint}
    \leanok
    \uses{def:qca_supported_in}
    If $X,Y\in\mathcal A_{\mathrm{loc}}$ are supported in finite regions
    $\Lambda$ and $\Gamma$ with $\Lambda\cap\Gamma=\varnothing$, then
    \begin{align}
        XY & =YX.
    \end{align}
    In particular, for $A\in\mathcal A_\Lambda$ and
    $B\in\mathcal A_\Gamma$,
    \begin{align}
        \iota_\Lambda(A)\iota_\Gamma(B)
         & =\iota_\Gamma(B)\iota_\Lambda(A).
    \end{align}
    This is the algebraic-local form of the tensor-factor identity underlying
    the construction in \cite[lines~2292--2300]{Cirac2017MPU}, not a named
    theorem of the appendix.
    % Source context: paper_v2.tex 2292--2300.
\end{lemma}

\begin{proof}\leanok
    \uses{thm:qca_disjoint_local_inclusion_commute}
    For representatives $X=\iota_\Lambda(A)$ and
    $Y=\iota_\Gamma(B)$, embedding both into
    $\mathcal A_{\Lambda\cup\Gamma}$ reduces the commutator to
    Lemma~\ref{thm:qca_disjoint_local_inclusion_commute} gives
    \begin{align}
        XY
         & =\iota_{\Lambda\cup\Gamma}
        \bigl(\iota_{\Lambda,\Lambda\cup\Gamma}(A)
        \iota_{\Gamma,\Lambda\cup\Gamma}(B)\bigr) \\
         & =\iota_{\Lambda\cup\Gamma}
        \bigl(\iota_{\Gamma,\Lambda\cup\Gamma}(B)
        \iota_{\Lambda,\Lambda\cup\Gamma}(A)\bigr)
        =YX.
    \end{align}
\end{proof}

\begin{lemma}[Algebraic-local relative commutant]
    \label{thm:qca_algebraic_local_relative_commutant}
    \lean{SpinChain.supportedIn_iff_commute_disjoint}
    \leanok
    \uses{def:qca_supported_in}
    Assume $d>0$. A local observable $X\in\mathcal A_{\mathrm{loc}}$ is
    supported in a finite region $\Lambda$ if and only if it commutes with every
    local observable having a finite support disjoint from $\Lambda$:
    \begin{align}
        X\in\iota_\Lambda(\mathcal A_\Lambda)
         & \iff
        \forall \Gamma,\ \forall Y,\quad
        \bigl(\Gamma\cap\Lambda=\varnothing\ \land\
        Y\in\iota_\Gamma(\mathcal A_\Gamma)\bigr)\Longrightarrow XY=YX.
    \end{align}
    This characterizes locally supported observables on the inductive limit
    algebra.
\end{lemma}

\begin{proof}\leanok
    \uses{thm:qca_finite_relative_commutant, thm:qca_disjoint_supported_local_commute, thm:qca_local_observable_finite_support}
    The forward implication is commutation of disjointly supported observables.
    Conversely, choose a finite support $\Delta$ of $X$ and represent $X$ by
    $A\in\mathcal A_{\Lambda\cup\Delta}$. The hypothesis says that $A$ commutes
    with every observable on $(\Lambda\cup\Delta)\setminus\Lambda$. Since
    $d>0$, the canonical map from $\mathcal A_{\Lambda\cup\Delta}$ into
    $\mathcal A_{\mathrm{loc}}$ is injective, so this commutation already holds
    in the finite-region algebra. The finite relative-commutant theorem gives
    $A=\iota_{\Lambda,\Lambda\cup\Delta}(A_\Lambda)$ for some
    $A_\Lambda\in\mathcal A_\Lambda$. Hence
    $X=\iota_\Lambda(A_\Lambda)$.
\end{proof}

\begin{lemma}[Quasi-local translation action laws]
    \label{lem:qca_quasi_local_translation_action}
    \lean{SpinChain.quasiLocalTranslation_zero,
        SpinChain.quasiLocalTranslation_add,
        SpinChain.quasiLocalTranslation_symm}
    \leanok
    \uses{def:qca_quasi_local_translation}
    For $d>0$, the quasi-local automorphisms form an additive action of
    $\mathbb Z$:
    \begin{align}
        \tau_0            & =\operatorname{id}, \\
        \tau_b\circ\tau_a & =\tau_{a+b},        \\
        \tau_a^{-1}       & =\tau_{-a}.
    \end{align}
\end{lemma}

\begin{proof}\leanok
    \uses{def:qca_quasi_local_translation, lem:qca_algebraic_local_translation_action}
    On the canonical image of the algebraic-local algebra,
    \begin{align}
        \tau_b\bigl(\tau_a(j(A))\bigr)
                     & =j\bigl(\tau_b(\tau_a(A))\bigr)
        =j(\tau_{a+b}(A)),                             \\
        \tau_0(j(A)) & =j(A).
    \end{align}
    Continuity extends these identities to the completion.  Taking
    $(a,b)=(a,-a)$ and $(-a,a)$ identifies the inverse of $\tau_a$ with
    $\tau_{-a}$.
\end{proof}

\begin{definition}[Translation covariance]
    \label{def:qca_translation_covariant}
    \lean{SpinChain.TranslationCovariant}
    \leanok
    \uses{def:qca_quasi_local_translation, def:qca_quasi_local_algebra}
    For $d>0$, a star-algebra automorphism $\omega$ of the quasi-local algebra
    $\mathcal A$ is translation covariant if
    \begin{align}
        \omega\circ\tau_a=\tau_a\circ\omega
    \end{align}
    for every $a\in\mathbb Z$.  In \cite[line~2298]{Cirac2017MPU}, this commutation relation is the translation-covariance condition in the
    definition of a one-dimensional QCA.
\end{definition}

\begin{lemma}[Pointwise characterization of translation covariance]
    \label{thm:qca_translation_covariant_iff}
    \lean{SpinChain.translationCovariant_iff}
    \leanok
    \uses{def:qca_translation_covariant}
    For $d>0$, a star-algebra automorphism $\omega$ of $\mathcal A$ is
    translation covariant if and only if
    \begin{align}
        \omega\bigl(\tau_a(A)\bigr)=\tau_a\bigl(\omega(A)\bigr)
    \end{align}
    for every $a\in\mathbb Z$ and every $A\in\mathcal A$.
\end{lemma}

\begin{proof}\leanok
    \uses{def:qca_translation_covariant}
    Two automorphisms are equal if and only if they agree on every observable,
    which applies directly to the two composites in the definition.
\end{proof}

\begin{lemma}[Translation covariance from unit translation]
    \label{thm:qca_translation_covariant_iff_commute_one}
    \lean{SpinChain.translationCovariant_iff_commute_one}
    \leanok
    \uses{def:qca_translation_covariant}
    For $d>0$, a star-algebra automorphism $\omega$ of $\mathcal A$ commutes with
    every integer translation if and only if it commutes with unit translation:
    \begin{align}
        \omega\circ\tau_1=\tau_1\circ\omega.
    \end{align}
\end{lemma}

\begin{proof}\leanok
    \uses{def:qca_translation_covariant, lem:qca_quasi_local_translation_action}
    The forward implication follows by taking the translation step to be $1$.
    Conversely, use two-sided integer induction.  The zero case follows from
    $\tau_0=\operatorname{id}$.  For the successor step, the additive action law
    gives $\tau_{i+1}=\tau_1\circ\tau_i$, and
    \begin{align}
        \omega\circ\tau_{i+1}
         & =\omega\circ\tau_1\circ\tau_i
        =\tau_1\circ\omega\circ\tau_i    \\
         & =\tau_1\circ\tau_i\circ\omega
        =\tau_{i+1}\circ\omega.
    \end{align}
    For the predecessor step, $\tau_{-1}=\tau_1^{-1}$ transfers commutation
    to $\tau_{-1}$, while $\tau_{i-1}=\tau_{-1}\circ\tau_i$, so
    \begin{align}
        \omega\circ\tau_{i-1}
         & =\omega\circ\tau_{-1}\circ\tau_i
        =\tau_{-1}\circ\omega\circ\tau_i    \\
         & =\tau_{-1}\circ\tau_i\circ\omega
        =\tau_{i-1}\circ\omega.
    \end{align}
\end{proof}

\begin{lemma}[Identity and composition of translation-covariant automorphisms]
    \label{thm:qca_translation_covariant_identity_composition}
    \lean{SpinChain.translationCovariant_refl,
        SpinChain.TranslationCovariant.trans}
    \leanok
    \uses{def:qca_translation_covariant}
    For $d>0$, the identity automorphism of $\mathcal A$ is translation
    covariant.  If $\omega$ and $\eta$ are translation covariant, then so is
    the composite
    \begin{align}
        A\longmapsto\eta\bigl(\omega(A)\bigr),
    \end{align}
    which first applies $\omega$ and then $\eta$.
\end{lemma}

\begin{proof}\leanok
    \uses{def:qca_translation_covariant}
    The identity commutes pointwise with every translation.  For the composite,
    translation covariance gives
    \begin{align}
        \eta\bigl(\omega(\tau_a(A))\bigr)
         & =\eta\bigl(\tau_a(\omega(A))\bigr)
        =\tau_a\bigl(\eta(\omega(A))\bigr).
    \end{align}
\end{proof}

\begin{lemma}[Norm isometry of quasi-local translation]
    \label{lem:qca_quasi_local_translation_isometry}
    \lean{SpinChain.norm_quasiLocalTranslation,
        SpinChain.quasiLocalTranslation_isometry}
    \leanok
    \uses{def:qca_quasi_local_translation, def:qca_quasi_local_algebra}
    For $d>0$, every quasi-local translation preserves the operator norm and
    distance:
    \begin{align}
        \lVert\tau_a(A)\rVert           & =\lVert A\rVert,   \\
        \lVert\tau_a(A)-\tau_a(B)\rVert & =\lVert A-B\rVert.
    \end{align}
\end{lemma}

\begin{proof}\leanok
    \uses{def:qca_quasi_local_translation, def:qca_quasi_local_algebra}
    A star-algebra equivalence between complex C-star algebras preserves the
    C-star norm.  Since $\tau_a$ is linear,
    \begin{align}
        \lVert\tau_a(A)-\tau_a(B)\rVert
         & =\lVert\tau_a(A-B)\rVert
        =\lVert A-B\rVert.
    \end{align}
\end{proof}

\begin{definition}[Finite-region support in the quasi-local algebra]
    \label{def:qca_quasi_local_supported_in}
    \lean{SpinChain.QuasiLocalSupportedIn,
        SpinChain.QuasiLocalSupportedIn.quasiLocalObservable,
        SpinChain.QuasiLocalSupportedIn.iff_exists}
    \leanok
    \uses{def:qca_quasi_local_observable}
    For $d>0$, a quasi-local observable $x\in\mathcal A$ is supported in a
    finite region $\Lambda\subset\mathbb Z$ if it belongs to the image of the
    canonical embedding $j_\Lambda\colon\mathcal A_\Lambda\to\mathcal A$.
    Equivalently,
    \begin{align}
        x\text{ is supported in }\Lambda
         & \iff x\in j_\Lambda(\mathcal A_\Lambda)         \\
         & \iff \text{there exists }A\in\mathcal A_\Lambda
        \text{ such that }j_\Lambda(A)=x.
    \end{align}
    In particular, $j_\Lambda(A)$ is supported in $\Lambda$ for every
    $A\in\mathcal A_\Lambda$.
\end{definition}

\begin{proof}\leanok
    \uses{def:qca_quasi_local_observable}
    These assertions follow directly from membership in the range of
    $j_\Lambda$.
\end{proof}

\begin{lemma}[Monotonicity of quasi-local support]
    \label{lem:qca_quasi_local_support_mono}
    \lean{SpinChain.QuasiLocalSupportedIn.mono}
    \leanok
    \uses{def:qca_quasi_local_supported_in}
    Let $d>0$.  If $x\in\mathcal A$ is supported in $\Lambda$ and
    $\Lambda\subseteq\Gamma$, then $x$ is supported in $\Gamma$.
\end{lemma}

\begin{proof}\leanok
    \uses{def:qca_quasi_local_supported_in, def:qca_local_inclusion, lem:qca_quasi_local_observable}
    For $x=j_\Lambda(A)$, enlarging $A$ to $\mathcal A_\Gamma$ by tensoring
    with the identity on $\Gamma\setminus\Lambda$ and using compatibility of the
    canonical embeddings identifies the image of this enlargement with $x$.
\end{proof}

\begin{lemma}[Disjointly supported quasi-local observables commute]
    \label{thm:qca_disjoint_quasi_local_commute}
    \lean{SpinChain.quasiLocalObservable_commute_of_disjoint,
        SpinChain.QuasiLocalSupportedIn.commute_of_disjoint}
    \leanok
    \uses{def:qca_quasi_local_supported_in}
    Let $d>0$.  If $x,y\in\mathcal A$ are supported in finite regions
    $\Lambda$ and $\Gamma$ with $\Lambda\cap\Gamma=\varnothing$, then
    \begin{align}
        xy & =yx.
    \end{align}
    Equivalently, for $A\in\mathcal A_\Lambda$ and
    $B\in\mathcal A_\Gamma$,
    \begin{align}
        j_\Lambda(A)j_\Gamma(B) & =j_\Gamma(B)j_\Lambda(A).
    \end{align}
    This is finite-support infrastructure for the quasi-local algebra in
    \cite[lines~2292--2300]{Cirac2017MPU}; it is not stated there as a separate
    theorem.
    % Source context: paper_v2.tex 2292--2300.
\end{lemma}

\begin{proof}\leanok
    \uses{thm:qca_disjoint_supported_local_commute}
    The completion map is a star-algebra homomorphism.  Applying it to the
    algebraic-local identity gives
    \begin{align}
        j(XY)=j(YX),
    \end{align}
    and finite-region support provides representatives
    $x=j_\Lambda(A)$ and $y=j_\Gamma(B)$.
\end{proof}

Commutation with observables on the complementary tensor factor characterizes
exact local support via the Weyl twirl.

\begin{definition}[Complementary-factor Weyl twirl]
    \label{def:qca_complementary_weyl_twirl}
    \lean{Matrix.reindexedWeyl,Matrix.complementaryWeylTwirl}
    \leanok
    Let $S$ and $C$ be finite sets, let $n=|C|>0$, and choose an equivalence
    $e\colon C\simeq\mathbb Z/n\mathbb Z$ and a primitive $n$-th root of unity
    $\zeta$.  If $W_{a,b}$ is the corresponding reindexed Weyl operator on
    $\mathbb C^C$, define
    \begin{align}
        \mathcal E_C(M)
         & = \frac{1}{n^2}\sum_{a,b\in\mathbb Z/n\mathbb Z}
        (\mathbf 1_S\otimes W_{a,b})M
        (\mathbf 1_S\otimes W_{a,b})^\dagger.
    \end{align}
\end{definition}

\begin{lemma}[Complementary-factor Weyl average]
    \label{lem:qca_complementary_weyl_twirl_eq}
    \lean{Matrix.sum_kronecker_one_reindexedWeyl_conj,
        Matrix.complementaryWeylTwirl_eq_partialTraceRight_kronecker,
        Matrix.complementaryWeylTwirl_eq_smul_partialTraceRight_kronecker_one}
    \leanok
    \uses{def:qca_complementary_weyl_twirl}
    For every $M\in\operatorname{Mat}_{S\times C}(\mathbb C)$,
    \begin{align}
        \mathcal E_C(M)
         & = \operatorname{tr}_C(M)\otimes\frac{\mathbf 1_C}{n}                \\
         & = \left(\frac{1}{n}\operatorname{tr}_C(M)\right)\otimes\mathbf 1_C.
    \end{align}
\end{lemma}

\begin{proof}\leanok
    Splitting $M$ into $C\times C$ blocks indexed by $S$ shows that the Weyl average sends
    each block $B$ to $(\operatorname{tr}B/n)\mathbf 1_C$, so its block trace is
    the corresponding entry of $\operatorname{tr}_C(M)$.
\end{proof}

\begin{lemma}[Norm estimate for the complementary-factor twirl]
    \label{lem:qca_complementary_weyl_twirl_norm}
    \lean{Matrix.norm_complementaryWeylTwirl_le,
        Matrix.norm_complementaryWeylTwirl_sub_le}
    \leanok
    \uses{def:qca_complementary_weyl_twirl}
    Let $S$ be a nonempty finite set.  The average is contractive in the
    operator norm; if every complementary Weyl conjugation fixes $X$, then
    \begin{align}
        \lVert\mathcal E_C(M)-X\rVert
         & \le \lVert M-X\rVert.
    \end{align}
\end{lemma}

\begin{proof}\leanok
    Unitary conjugation preserves the operator norm, and the coefficients of the
    finite average sum to one; applying this contraction to $M-X$ gives the bound.
\end{proof}

\begin{lemma}[Exact support from outside commutation]
    \label{thm:qca_quasi_local_supported_iff_commute_disjoint}
    \lean{SpinChain.quasiLocalSupportedIn_iff_commute_disjoint}
    \leanok
    \uses{def:qca_quasi_local_supported_in}
    Let $d>0$, let $\Lambda\subset\mathbb Z$ be finite, and let
    $x\in\mathcal A$.  Then
    \begin{align}
        x\in\mathcal A_\Lambda
        \quad\Longleftrightarrow\quad
        [x,y]=0
        \text{ for every }y\in\mathcal A_\Gamma
        \text{ and every finite }\Gamma\cap\Lambda=\varnothing.
    \end{align}
    As in \cite{Cirac2017MPU,SchumacherWerner2004QCA}, this equivalence characterizes local support through commutants on complementary
    regions in the UHF algebra.
\end{lemma}

\begin{proof}\leanok
    \uses{thm:qca_disjoint_quasi_local_commute, lem:qca_quasi_local_observable, thm:qca_quasi_local_density, lem:qca_local_inclusion_norm, thm:mpu_reindex_unitary, def:qca_complementary_weyl_twirl, lem:qca_complementary_weyl_twirl_eq}
    The forward implication follows from commutation of disjoint tensor factors.
    Conversely, approximating $x$ by $A\in\mathcal A_\Delta$ with
    $\Lambda\subseteq\Delta$, splitting
    $\Delta=\Lambda\sqcup(\Delta\setminus\Lambda)$, and averaging $A$ over Weyl
    unitaries on the second factor produces an approximant in $\mathcal A_\Lambda$: every such conjugation fixes $x$ by the
    commutation hypothesis, so unitary invariance and convexity give
    \begin{align}
        \lVert\mathcal E_{\Delta\setminus\Lambda}(A)-x\rVert
         & \le \lVert A-x\rVert.
    \end{align}
    The Weyl identity writes the average as $B\otimes\mathbf 1$ for some
    $B\in\mathcal A_\Lambda$.  Thus elements of $\mathcal A_\Lambda$ approximate
    $x$ arbitrarily well.  The embedding of $\mathcal A_\Lambda$ is isometric,
    hence has closed range, and therefore $x\in\mathcal A_\Lambda$.
\end{proof}

As used in \cite{Cirac2017MPU,Gross2012IndexTheory,SchumacherWerner2004QCA}, the homogeneous quasi-local UHF algebra has a trivial center consisting of
scalar multiples of the
identity.

\begin{lemma}[Scalar center of the homogeneous quasi-local algebra]
    \label{thm:qca_quasi_local_center_scalar}
    \lean{SpinChain.existsUnique_eq_smul_one_of_commute_quasiLocalObservable,
        SpinChain.existsUnique_eq_smul_one_of_commute,
        SpinChain.instIsCentralQuasiLocalAlgebra,
        SpinChain.quasiLocalAlgebra_center_eq_bot}
    \leanok
    \uses{def:qca_quasi_local_algebra, def:qca_quasi_local_observable}
    Let $d>0$ and let $z\in\mathcal A$.  If
    \begin{align}
        \forall\Gamma\subset\mathbb Z\text{ finite},\quad
        \forall A\in\mathcal A_\Gamma,\qquad
        z j_\Gamma(A) & =j_\Gamma(A)z,
    \end{align}
    then there is a unique $c\in\mathbb C$ such that
    \begin{align}
        z & =c\mathbf 1.
    \end{align}
    Consequently,
    \begin{align}
        \{z\in\mathcal A\mid za=az\text{ for every }a\in\mathcal A\}
         & =\mathbb C\mathbf 1.
    \end{align}
\end{lemma}

\begin{proof}\leanok
    \uses{thm:qca_quasi_local_supported_iff_commute_disjoint, def:qca_quasi_local_supported_in, lem:qca_quasi_local_observable}
    By Lemma~\ref{thm:qca_quasi_local_supported_iff_commute_disjoint}
    with $\Lambda=\varnothing$, the commutation hypothesis gives
    $z\in j_\varnothing(\mathcal A_\varnothing)$, so there is an
    $A\in\mathcal A_\varnothing$ with $z=j_\varnothing(A)$.  Let
    $e\colon\mathbb C\xrightarrow{\ \sim\ }\mathcal A_\varnothing$ be the
    canonical isomorphism and write $A=e(c)$.  If also
    $z=j_\varnothing(e(c'))$, then injectivity of $j_\varnothing$ and of $e$
    give
    \begin{align}
        j_\varnothing(e(c))=j_\varnothing(e(c'))
         & \Longrightarrow e(c)=e(c')
        \Longrightarrow c=c'.
    \end{align}
    Thus $c$ is the unique scalar for which $z=c\mathbf 1$.

    With $Z(\mathcal A)$ denoting the center, since every
    $j_\Gamma(B)$ belongs to $\mathcal A$, the two inclusions are
    \begin{align}
        z\in Z(\mathcal A)
         & \Longrightarrow
        \bigl(\forall \Gamma\subset\mathbb Z\text{ finite},\
        \forall B\in\mathcal A_\Gamma,\
        z j_\Gamma(B)=j_\Gamma(B)z\bigr)
        \Longrightarrow z\in\mathbb C\mathbf 1, \\
        z=c\mathbf 1\quad(c\in\mathbb C)
         & \Longrightarrow
        \bigl(\forall a\in\mathcal A,\
        (c\mathbf 1)a=a(c\mathbf 1)\bigr)
        \Longrightarrow z\in Z(\mathcal A).
    \end{align}
    Hence $Z(\mathcal A)=\mathbb C\mathbf 1$.
\end{proof}

\section{Quantum cellular automata}

The appendix requires that ``there exists a finite subset
$\mathcal N\subset\mathbb Z$'' such that, for every finite
$\Lambda\subset\mathbb Z$,
\begin{align}
    \omega(\mathcal A_\Lambda) & \subset
    \mathcal A_{\Lambda+\mathcal N}.
\end{align}
In \cite[line~2298]{Cirac2017MPU}, this inclusion is the forward propagation
condition.

\begin{definition}[Propagation within a fixed neighborhood]
    \label{def:qca_propagates_within}
    \lean{SpinChain.PropagatesWithin}
    \leanok
    \uses{def:qca_quasi_local_supported_in, def:qca_region_sumset}
    Let $d>0$, let $\omega\colon\mathcal A\to\mathcal A$ be a star-algebra
    automorphism, and let $\mathcal N\subset\mathbb Z$ be finite.  The
    automorphism $\omega$ propagates within $\mathcal N$ if, for every finite
    region $\Lambda$ and every $x\in\mathcal A$,
    \begin{align}
        x\text{ supported in }\Lambda
         & \implies
        \omega(x)\text{ supported in }\Lambda+\mathcal N.
    \end{align}
\end{definition}

\begin{definition}[Finite forward propagation]
    \label{def:qca_finite_propagation}
    \lean{SpinChain.HasFinitePropagation}
    \leanok
    \uses{def:qca_propagates_within}
    Let $d>0$ and let $\omega\colon\mathcal A\to\mathcal A$ be a star-algebra
    automorphism.  The automorphism $\omega$ has finite forward propagation if
    there exists a finite set $\mathcal N\subset\mathbb Z$ such that $\omega$
    propagates within $\mathcal N$.
\end{definition}

\begin{lemma}[Equivalent formulations of forward propagation]
    \label{lem:qca_propagates_within_iff}
    \lean{SpinChain.propagatesWithin_iff_quasiLocalObservable,
        SpinChain.propagatesWithin_iff_range_subset,
        SpinChain.propagatesWithin_iff_exists_local}
    \leanok
    \uses{def:qca_quasi_local_supported_in, def:qca_propagates_within, def:qca_quasi_local_observable, def:qca_region_sumset}
    Let $d>0$, let $\omega\colon\mathcal A\to\mathcal A$ be a star-algebra
    automorphism, and let $\mathcal N\subset\mathbb Z$ be finite.  The following
    conditions are equivalent:
    \begin{align}
         & \omega\text{ propagates within }\mathcal N; \tag{a}              \\
         & \forall\Lambda\Subset\mathbb Z,\ \forall A\in\mathcal A_\Lambda,
        \quad \omega(j_\Lambda(A))\text{ is supported in }
        \Lambda+\mathcal N; \tag{b}                                         \\
         & \forall\Lambda\Subset\mathbb Z,\quad
        \omega\bigl(j_\Lambda(\mathcal A_\Lambda)\bigr)
        \subseteq j_{\Lambda+\mathcal N}
        (\mathcal A_{\Lambda+\mathcal N}); \tag{c}                          \\
         & \forall\Lambda\Subset\mathbb Z,\ \forall A\in\mathcal A_\Lambda,
        \ \exists B\in\mathcal A_{\Lambda+\mathcal N},\quad
        \omega(j_\Lambda(A))=j_{\Lambda+\mathcal N}(B). \tag{d}
    \end{align}
    Thus condition~\textup{(c)} is precisely the range-inclusion formula from
    line~2298 of \cite{Cirac2017MPU}, expressed inside the quasi-local algebra,
    and condition~\textup{(d)} is its finite-region approximation form.
\end{lemma}

\begin{proof}\leanok
    \uses{def:qca_quasi_local_supported_in, def:qca_propagates_within, def:qca_quasi_local_observable}
    Every observable supported in $\Lambda$ has the form $j_\Lambda(A)$, so
    \textup{(a)} and \textup{(b)} are equivalent.  Replacing support by
    membership in the range of the canonical embedding gives \textup{(c)}.
    Finally, membership in
    $j_{\Lambda+\mathcal N}(\mathcal A_{\Lambda+\mathcal N})$ is equivalent to
    the existence of a representative $B$ in that finite-region algebra, which
    gives \textup{(d)}.
\end{proof}

\begin{lemma}[Enlargement of propagation neighborhoods]
    \label{lem:qca_propagation_neighborhood_mono}
    \lean{SpinChain.PropagatesWithin.mono,
        SpinChain.HasFinitePropagation.exists_superset}
    \leanok
    \uses{def:qca_propagates_within, def:qca_finite_propagation}
    Let $d>0$ and let $\omega\colon\mathcal A\to\mathcal A$ be a star-algebra
    automorphism.  If $\mathcal N\subseteq\mathcal M$ and $\omega$ propagates
    within $\mathcal N$, then $\omega$ propagates within $\mathcal M$.  More
    explicitly, for every finite region $\Lambda$,
    \begin{align}
        \Lambda+\mathcal N & \subseteq\Lambda+\mathcal M,                              \\
        x\text{ supported in }\Lambda
                           & \implies \omega(x)\text{ supported in }\Lambda+\mathcal N
        \implies \omega(x)\text{ supported in }\Lambda+\mathcal M.
    \end{align}
    Consequently, if $\omega$ has finite forward propagation, then for every
    prescribed finite set $\mathcal N\subset\mathbb Z$ there is a finite set
    $\mathcal M\subset\mathbb Z$ such that
    \begin{align}
        \mathcal N & \subseteq\mathcal M,                 \\
                   & \forall\Lambda\Subset\mathbb Z,\quad
        \omega(\mathcal A_\Lambda)\subseteq
        \mathcal A_{\Lambda+\mathcal M}.
    \end{align}
    For the propagation condition at Appendix line~2298 of \cite{Cirac2017MPU}, this monotonicity is the neighborhood-enlargement
    consequence.
\end{lemma}

\begin{proof}\leanok
    \uses{def:qca_region_sumset, lem:qca_quasi_local_support_mono, def:qca_propagates_within, def:qca_finite_propagation}
    Monotonicity of the sumset in its second argument gives
    $\Lambda+\mathcal N\subseteq\Lambda+\mathcal M$.  Monotonicity of
    quasi-local support then enlarges the support of $\omega(x)$ from
    $\Lambda+\mathcal N$ to $\Lambda+\mathcal M$.  For the final assertion,
    choose a finite propagation neighborhood $\mathcal K$ and take
    $\mathcal M=\mathcal K\cup\mathcal N$.
\end{proof}

\begin{lemma}[Composition of forward propagation bounds]
    \label{lem:qca_forward_propagation_comp}
    \lean{SpinChain.PropagatesWithin.trans,
        SpinChain.HasFinitePropagation.trans}
    \leanok
    \uses{def:qca_region_sumset, def:qca_propagates_within, def:qca_finite_propagation}
    Let $d>0$, let $\omega,\eta\colon\mathcal A\to\mathcal A$ be star-algebra
    automorphisms, and let $\mathcal N,\mathcal M\subset\mathbb Z$ be finite.
    Suppose that $\omega$ propagates within $\mathcal N$ and $\eta$ propagates
    within $\mathcal M$.  The composite $\eta\circ\omega$ first applies
    $\omega$ and then $\eta$, so
    \begin{align}
        (\eta\circ\omega)(x) & =\eta(\omega(x)),                                         \\
        x\text{ supported in }\Lambda
                             & \implies \omega(x)\text{ supported in }\Lambda+\mathcal N \\
                             & \implies \eta(\omega(x))\text{ supported in }
        (\Lambda+\mathcal N)+\mathcal M
        =\Lambda+(\mathcal N+\mathcal M).
    \end{align}
    Hence $\eta\circ\omega$ propagates within $\mathcal N+\mathcal M$.
    In particular, if $\omega$ and $\eta$ have finite forward propagation, then
    $\eta\circ\omega$ has finite forward propagation, establishing for Appendix line~2298 of \cite{Cirac2017MPU} the composition
    closure of the finite-neighborhood
    condition.
\end{lemma}

\begin{proof}\leanok
    \uses{def:qca_region_sumset, def:qca_propagates_within, def:qca_finite_propagation}
    Successive application of the two propagation bounds and associativity of finite
    region sumsets give the containment; for the existential assertion, choosing one finite propagation
    neighborhood for each automorphism and taking their sumset completes the proof.
\end{proof}

\begin{lemma}[Symmetric-interval normalization of finite propagation]
    \label{lem:qca_finite_propagation_symmetric_interval}
    \lean{SpinChain.exists_subset_symmetric_Icc,
        SpinChain.PropagatesWithin.exists_symmetric_Icc,
        SpinChain.HasFinitePropagation.exists_symmetric_Icc}
    \leanok
    \uses{def:qca_propagates_within, def:qca_finite_propagation}
    Every finite set $\mathcal N\subset\mathbb Z$ is contained in
    $[-R,R]\cap\mathbb Z$ for some $R\in\mathbb N$.  This includes
    $\mathcal N=\varnothing$, for which one may take $R=0$.  If $\omega$
    propagates within a fixed set $\mathcal N$, then $R$ may be chosen so that
    \begin{align}
        \mathcal N & \subseteq[-R,R]\cap\mathbb Z,                  \\
        \omega     & \text{ propagates within }[-R,R]\cap\mathbb Z.
    \end{align}
    Consequently, finite forward propagation admits a symmetric-interval
    bound:
    \begin{align}
        \exists R\in\mathbb N,\qquad
        \omega & \text{ propagates within }[-R,R]\cap\mathbb Z.
    \end{align}
\end{lemma}

\begin{proof}\leanok
    \uses{def:qca_finite_propagation, lem:qca_propagation_neighborhood_mono}
    Let $R$ be the maximum of $|n|$ over $n\in\mathcal N$, with value $0$
    when $\mathcal N$ is empty.  Then $-R\leq n\leq R$ for every
    $n\in\mathcal N$, which gives
    $\mathcal N\subseteq[-R,R]\cap\mathbb Z$.  Enlargement of propagation
    neighborhoods gives the fixed-neighborhood assertion.  Finally, choose a
    finite propagation neighborhood and apply this enlargement to its symmetric
    interval.
\end{proof}

\begin{lemma}[Disjointness for reflected neighborhoods]
    \label{lem:qca_disjoint_region_sumset_neg_iff}
    \lean{SpinChain.disjoint_regionSumset_neg_iff}
    \leanok
    \uses{def:qca_region_sumset}
    Let $\Gamma,\Lambda,\mathcal N\subset\mathbb Z$ be finite.  Then
    \begin{align}
        \Gamma\cap(\Lambda+(-\mathcal N))=\varnothing
        \quad\Longleftrightarrow\quad
        (\Gamma+\mathcal N)\cap\Lambda=\varnothing.
    \end{align}
\end{lemma}

\begin{proof}\leanok
    If $g+n=\lambda$ for $g\in\Gamma$, $n\in\mathcal N$, and
    $\lambda\in\Lambda$, then
    $g=\lambda+(-n)\in\Lambda+(-\mathcal N)$, proving the forward implication.
    Conversely, if $g\in\Gamma$ and $g=\lambda+(-n)$ with
    $\lambda\in\Lambda$ and $n\in\mathcal N$, then
    $\lambda=g+n\in\Gamma+\mathcal N$.
\end{proof}

\begin{lemma}[Inverse propagation and inverse covariance]
    \label{thm:qca_inverse_propagation}
    \lean{SpinChain.PropagatesWithin.symm,
        SpinChain.HasFinitePropagation.symm,
        SpinChain.TranslationCovariant.symm}
    \leanok
    \uses{def:qca_propagates_within, def:qca_finite_propagation, def:qca_translation_covariant}
    Let $d>0$, let $\omega$ be a star-algebra automorphism of $\mathcal A$, and
    let $\mathcal N\subset\mathbb Z$ be finite.  If $\omega$ propagates within
    $\mathcal N$, then $\omega^{-1}$ propagates within $-\mathcal N$.  Therefore
    finite forward propagation of $\omega$ implies finite forward propagation of
    $\omega^{-1}$.  If $\omega$ is translation covariant, then $\omega^{-1}$
    is also translation covariant.

    Schumacher and Werner \cite[Theorem~6, Corollary~7]{SchumacherWerner2004QCA}
    derive inverse locality through generalized Margolus partitioning; the
    proof here establishes it via commutants on complementary regions.
\end{lemma}

\begin{proof}\leanok
    \uses{lem:qca_disjoint_region_sumset_neg_iff, thm:qca_quasi_local_supported_iff_commute_disjoint}
    Let $x\in\mathcal A_\Lambda$ and let
    $y\in\mathcal A_\Gamma$, where
    $\Gamma\cap(\Lambda+(-\mathcal N))=\varnothing$.  Forward locality gives
    $\omega(y)\in\mathcal A_{\Gamma+\mathcal N}$.  The regions
    $\Gamma+\mathcal N$ and $\Lambda$ are disjoint, so $x$ commutes with
    $\omega(y)$.  Applying $\omega^{-1}$ shows that $\omega^{-1}(x)$ commutes
    with $y$.  The exact outside-commutant characterization now gives
    \begin{align}
        \omega^{-1}(x) & \in\mathcal A_{\Lambda+(-\mathcal N)}.
    \end{align}
    Choosing the reflected finite neighborhood proves the existential statement.
    For each lattice translation $\tau_a$, multiplying the covariance identity
    on the left and right by $\omega^{-1}$ gives
    \begin{align}
        \omega\circ\tau_a=\tau_a\circ\omega
         & \implies
        \omega^{-1}\circ\tau_a=\tau_a\circ\omega^{-1}.
    \end{align}
\end{proof}

\begin{definition}[One-dimensional quantum cellular automaton]
    \label{def:qca}
    \lean{SpinChain.IsQCA}
    \leanok
    \uses{def:qca_translation_covariant, def:qca_finite_propagation}
    Let $d>0$ and let $\omega\colon\mathcal A\to\mathcal A$ be a star-algebra
    automorphism.  The automorphism $\omega$ is a one-dimensional quantum
    cellular automaton if it is translation covariant and has finite forward
    propagation.  Thus there exists a finite set $\mathcal N\subset\mathbb Z$
    such that, for every finite $\Lambda\subset\mathbb Z$,
    \begin{align}
        \omega(\mathcal A_\Lambda) & \subseteq
        \mathcal A_{\Lambda+\mathcal N}.
    \end{align}
    At line~2298 of \cite{Cirac2017MPU}, this definition imposes only
    the displayed forward propagation condition and makes no locality claim for
    $\omega^{-1}$.
\end{definition}

\begin{lemma}[Exact interface for quantum cellular automata]
    \label{thm:qca_iff_mk_projections}
    \lean{SpinChain.isQCA_iff,SpinChain.IsQCA.mk,
        SpinChain.IsQCA.translationCovariant,
        SpinChain.IsQCA.hasFinitePropagation}
    \leanok
    \uses{def:qca}
    For $d>0$, a star-algebra automorphism $\omega$ of $\mathcal A$ is a
    quantum cellular automaton if and only if it is translation covariant and
    has finite forward propagation.  These two properties therefore construct an
    automaton, and every quantum cellular automaton satisfies each of them.
\end{lemma}

\begin{proof}\leanok
    \uses{def:qca}
    This is the conjunction defining a quantum cellular automaton, together with
    its introduction and elimination rules.
\end{proof}

\begin{lemma}[Identity quantum cellular automaton]
    \label{thm:qca_identity}
    \lean{SpinChain.isQCA_refl}
    \leanok
    \uses{def:qca}
    For $d>0$, the identity automorphism of $\mathcal A$ is a quantum cellular
    automaton.
\end{lemma}

\begin{proof}\leanok
    \uses{def:qca_finite_propagation, def:qca_propagates_within, def:qca_region_sumset, thm:qca_iff_mk_projections, thm:qca_translation_covariant_identity_composition}
    The identity is translation covariant.  For finite forward propagation,
    choose the neighborhood $\{0\}$.  An observable supported in $\Lambda$
    remains unchanged, and $\Lambda+\{0\}=\Lambda$.
\end{proof}

\begin{lemma}[Composition of quantum cellular automata]
    \label{thm:qca_composition}
    \lean{SpinChain.IsQCA.trans}
    \leanok
    \uses{def:qca}
    Let $d>0$.  If $\omega$ and $\eta$ are quantum cellular automata on
    $\mathcal A$, then $\eta\circ\omega$ is a quantum cellular automaton.  This
    composite first applies $\omega$ and then $\eta$.  No propagation condition
    for either inverse is required.
\end{lemma}

\begin{proof}\leanok
    \uses{thm:qca_iff_mk_projections, thm:qca_translation_covariant_identity_composition, lem:qca_forward_propagation_comp}
    Both defining properties are closed under composition in the stated order.
    Translation covariance of $\omega$ and $\eta$ gives translation covariance
    of $\eta\circ\omega$.  If $\mathcal N$ and $\mathcal M$ are forward
    propagation neighborhoods for $\omega$ and $\eta$, respectively, then
    $\mathcal N+\mathcal M$ is a forward propagation neighborhood for
    $\eta\circ\omega$.  Hence the composite is a quantum cellular automaton.
\end{proof}

\begin{lemma}[Inverse quantum cellular automaton]
    \label{thm:qca_inverse}
    \lean{SpinChain.IsQCA.symm}
    \leanok
    \uses{def:qca}
    Let $d>0$.  If $\omega$ is a quantum cellular automaton on $\mathcal A$, then
    $\omega^{-1}$ is also a quantum cellular automaton.
\end{lemma}

\begin{proof}\leanok
    \uses{thm:qca_iff_mk_projections, thm:qca_inverse_propagation}
    Translation covariance and finite forward propagation are both preserved by
    inversion.  These are the two defining properties of a quantum cellular
    automaton.
\end{proof}

\section{Grouping the sites of a quantum cellular automaton}

\begin{lemma}[Support transformation under site blocking]
    \label{lem:qca_quasi_local_blocking_support}
    \lean{SpinChain.quasiLocalBlocking_supportedIn,
        SpinChain.quasiLocalBlocking_symm_supportedIn,
        SpinChain.quasiLocalBlocking_supportedIn_expandedRegion_iff}
    \leanok
    \uses{lem:qca_quasi_local_blocking_formulas, def:qca_quasi_local_supported_in, lem:qca_block_hull_properties}
    Let $d>0$ and $L>0$.  If a blocked-chain observable $X$ is supported in
    $\Lambda$, then $\overline{\mathcal B}_L(X)$ is supported in
    $\widehat\Lambda_L$.  If an original-chain observable $Y$ is supported in
    $\Delta$, then $\overline{\mathcal B}_L^{-1}(Y)$ is supported in
    $\operatorname{Hull}_L(\Delta)$.  In particular, support preservation is exact
    on expanded regions:
    \begin{align}
        \overline{\mathcal B}_L(X)
        \text{ is supported in }\widehat\Lambda_L
        \quad\Longleftrightarrow\quad
        X\text{ is supported in }\Lambda.
    \end{align}
\end{lemma}

\begin{proof}\leanok
    \uses{lem:qca_quasi_local_blocking_formulas, lem:qca_block_hull_properties}
    Applying the corresponding forward or inverse finite evaluation formula to a finite-region representative
    gives the forward support bounds.  For the reverse implication in the
    exact statement, apply inverse blocking and simplify the inverse composite
    and the block hull of the expanded region.
\end{proof}

In the representation theorem of \cite[line~2308]{Cirac2017MPU}, finite blocking preserves the one-dimensional quantum cellular automaton
structure and provides the blocked site
grouping.

\begin{definition}[Blocked automorphism]
    \label{def:qca_blocked_automorphism}
    \lean{SpinChain.blockedAutomorphism}
    \leanok
    \uses{def:qca_quasi_local_blocking}
    Let $d,L>0$, let
    $\overline{\mathcal B}_L\colon\mathcal A^{(d^L)}\simeq\mathcal A^{(d)}$
    be the quasi-local blocking equivalence, and let $\omega$ be a star-algebra
    automorphism of $\mathcal A^{(d)}$.  The induced automorphism of the blocked
    chain is
    \begin{align}
        \omega^{[L]}
         & =\overline{\mathcal B}_L^{-1}\circ\omega\circ
        \overline{\mathcal B}_L.
    \end{align}
\end{definition}

\begin{lemma}[Evaluation of the blocked automorphism]
    \label{lem:qca_blocked_automorphism_apply}
    \lean{SpinChain.blockedAutomorphism_apply}
    \leanok
    \uses{def:qca_blocked_automorphism}
    For every $X\in\mathcal A^{(d^L)}$,
    \begin{align}
        \omega^{[L]}(X)
         & =\overline{\mathcal B}_L^{-1}
        \bigl(\omega(\overline{\mathcal B}_L(X))\bigr).
    \end{align}
\end{lemma}

\begin{proof}\leanok
    \uses{def:qca_blocked_automorphism}
    This is the evaluation formula for the defining conjugation.
\end{proof}

\begin{lemma}[Inverse of the blocked automorphism]
    \label{lem:qca_blocked_automorphism_symm}
    \lean{SpinChain.blockedAutomorphism_symm}
    \leanok
    \uses{def:qca_blocked_automorphism}
    The inverse of the blocked automorphism is the blocked inverse:
    \begin{align}
        (\omega^{[L]})^{-1}
         & =(\omega^{-1})^{[L]}.
    \end{align}
\end{lemma}

\begin{proof}\leanok
    \uses{lem:qca_blocked_automorphism_apply}
    On any blocked-chain observable, the blocking equivalence
    and its inverse cancel on the two sides of $\omega^{-1}$.
\end{proof}

\begin{lemma}[Propagation neighborhood after site blocking]
    \label{lem:qca_propagates_within_blocked}
    \lean{SpinChain.PropagatesWithin.blocked}
    \leanok
    \uses{def:qca_blocked_automorphism, def:qca_blocked_neighborhood, def:qca_propagates_within}
    Let $\mathcal N\subset\mathbb Z$ be finite.  If $\omega$ propagates within
    $\mathcal N$, then $\omega^{[L]}$ propagates within the carry-aware blocked
    neighborhood $\mathcal C_L(\mathcal N)$.  Explicitly, for every finite
    blocked region $\Lambda$ and every $X\in\mathcal A^{(d^L)}$ supported in
    $\Lambda$, the observable $\omega^{[L]}(X)$ is supported in
    $\Lambda+\mathcal C_L(\mathcal N)$.
\end{lemma}

\begin{proof}\leanok
    \uses{lem:qca_blocked_automorphism_apply, lem:qca_quasi_local_blocking_support, lem:qca_block_hull_region_sumset}
    Forward blocking sends support in $\Lambda$ to support in
    $\widehat\Lambda_L$, and applying the propagation bound for $\omega$ followed by
    inverse blocking places the image in
    \begin{align}
        \operatorname{Hull}_L
        \bigl(\widehat\Lambda_L+\mathcal N\bigr)
         & =\Lambda+\mathcal C_L(\mathcal N)
    \end{align}
    by the exact carry-aware block-hull identity.
\end{proof}

\begin{lemma}[Finite propagation after site blocking]
    \label{lem:qca_finite_propagation_blocked}
    \lean{SpinChain.HasFinitePropagation.blocked}
    \leanok
    \uses{def:qca_blocked_automorphism, def:qca_finite_propagation, lem:qca_propagates_within_blocked}
    If $\omega$ has finite forward propagation, then $\omega^{[L]}$ has finite
    forward propagation.
\end{lemma}

\begin{proof}\leanok
    \uses{lem:qca_propagates_within_blocked}
    For any finite propagation neighborhood $\mathcal N$ of $\omega$,
    Lemma~\ref{lem:qca_propagates_within_blocked} gives the finite neighborhood
    $\mathcal C_L(\mathcal N)$ for $\omega^{[L]}$.
\end{proof}

\begin{lemma}[Translation covariance of the blocked automorphism]
    \label{thm:qca_blocked_translation_covariant_iff}
    \lean{SpinChain.translationCovariant_blockedAutomorphism_iff}
    \leanok
    \uses{def:qca_blocked_automorphism, def:qca_translation_covariant, def:qca_quasi_local_translation}
    The blocked automorphism is translation covariant if and only if the
    original automorphism commutes with translation by one whole block:
    \begin{align}
        \omega^{[L]}\text{ is translation covariant}
        \quad\Longleftrightarrow\quad
        \omega\circ\overline\tau^{(d)}_L
         & =\overline\tau^{(d)}_L\circ\omega.
    \end{align}
    Thus blocked unit-translation covariance records commutation of $\omega$
    with translation by $L$, not necessarily commutation with original unit
    translation.  For $L\neq1$, the latter condition can be strictly stronger.
\end{lemma}

\begin{proof}\leanok
    \uses{lem:qca_blocked_automorphism_apply, thm:qca_translation_covariant_iff_commute_one, lem:qca_quasi_local_blocking_translation}
    By Lemma~\ref{thm:qca_translation_covariant_iff_commute_one}, translation
    covariance on the blocked chain is equivalent to commutation with blocked
    unit translation.  Under $\overline{\mathcal B}_L$, blocked unit translation
    corresponds to original translation by $L$.  Conjugating the commutation
    identity by the blocking equivalence proves both implications.
\end{proof}

\begin{lemma}[Translation covariance after site blocking]
    \label{lem:qca_translation_covariant_blocked}
    \lean{SpinChain.TranslationCovariant.blocked}
    \leanok
    \uses{def:qca_blocked_automorphism, def:qca_translation_covariant, thm:qca_blocked_translation_covariant_iff}
    If $\omega$ is translation covariant, then $\omega^{[L]}$ is translation
    covariant.
\end{lemma}

\begin{proof}\leanok
    \uses{thm:qca_blocked_translation_covariant_iff}
    Translation covariance of $\omega$ gives commutation with
    $\overline\tau^{(d)}_L$, so
    Lemma~\ref{thm:qca_blocked_translation_covariant_iff} gives translation covariance of $\omega^{[L]}$.
\end{proof}

\begin{lemma}[Blocked quantum cellular automaton]
    \label{lem:qca_blocked}
    \lean{SpinChain.IsQCA.blocked}
    \leanok
    \uses{def:qca, def:qca_blocked_automorphism, lem:qca_finite_propagation_blocked, lem:qca_translation_covariant_blocked}
    Let $d,L>0$.  If $\omega$ is a one-dimensional QCA on
    $\mathcal A^{(d)}$, then
    \begin{align}
        \omega^{[L]}
         & =\overline{\mathcal B}_L^{-1}\circ\omega\circ
        \overline{\mathcal B}_L
    \end{align}
    is a one-dimensional QCA on $\mathcal A^{(d^L)}$.
\end{lemma}

\begin{proof}\leanok
    \uses{lem:qca_finite_propagation_blocked, lem:qca_translation_covariant_blocked}
    The two defining properties of a one-dimensional QCA are finite forward
    propagation and translation covariance.  They are preserved by
    Lemma~\ref{lem:qca_finite_propagation_blocked} and
    Lemma~\ref{lem:qca_translation_covariant_blocked}, respectively.
\end{proof}

As in \cite[line~2308]{Cirac2017MPU} and \cite{SchumacherWerner2004QCA}, regrouping sites normalizes the finite propagation neighborhood of a QCA to
$[-1,1]\cap\mathbb Z$ and sets up the depth-two circuit
representation.

\begin{lemma}[Nearest-neighbor carries from a strict displacement bound]
    \label{lem:qca_blocked_neighborhood_subset_nearest_neighbor}
    \lean{SpinChain.blockedNeighborhood_subset_Icc_one}
    \leanok
    \uses{def:qca_blocked_neighborhood}
    Let $L>0$ and let $\mathcal N\subset\mathbb Z$ be finite.  If
    $|n|<L$ for every $n\in\mathcal N$, then
    \begin{align}
        \mathcal C_L(\mathcal N) & \subseteq[-1,1]\cap\mathbb Z.
    \end{align}
\end{lemma}

\begin{proof}\leanok
    \uses{lem:qca_blocked_neighborhood_membership}
    If $\mathcal N=\varnothing$, membership in
    $\mathcal C_L(\mathcal N)$ is impossible.  In general, a point
    $b\in\mathcal C_L(\mathcal N)$ yields residues $0\leq r,s<L$ and
    $n\in\mathcal N$ such that
    \begin{align}
        r+n & =bL+s.
    \end{align}
    Since $-L<n<L$, the equality is incompatible with $b\leq-2$ and with
    $b\geq2$.  Thus $-1\leq b\leq1$, including for negative displacements.
\end{proof}

\begin{lemma}[Supremum criterion for nearest-neighbor carries]
    \label{lem:qca_blocked_neighborhood_subset_nearest_neighbor_sup}
    \lean{SpinChain.blockedNeighborhood_subset_Icc_one_of_sup}
    \leanok
    \uses{def:qca_blocked_neighborhood}
    Let $L>0$ and let $\mathcal N\subset\mathbb Z$ be finite.  If
    \begin{align}
        \sup_{n\in\mathcal N}|n| & <L,
    \end{align}
    then
    \begin{align}
        \mathcal C_L(\mathcal N) & \subseteq[-1,1]\cap\mathbb Z.
    \end{align}
\end{lemma}

\begin{proof}\leanok
    \uses{lem:qca_blocked_neighborhood_subset_nearest_neighbor}
    Every $|n|$ is bounded above by the finite supremum, so the strict
    supremum bound gives the pointwise hypothesis.  When
    $\mathcal N=\varnothing$, the supremum is $0$ and the conclusion is
    vacuous.
\end{proof}

\begin{lemma}[Canonical nearest-neighbor block length]
    \label{lem:qca_blocked_neighborhood_canonical_nearest_neighbor}
    \lean{SpinChain.blockedNeighborhood_subset_Icc_one_sup_add_one}
    \leanok
    \uses{def:qca_blocked_neighborhood}
    For every finite $\mathcal N\subset\mathbb Z$, define
    \begin{align}
        L_{\mathcal N} & =1+\sup_{n\in\mathcal N}|n|.
    \end{align}
    Then $L_{\mathcal N}>0$ and
    \begin{align}
        \mathcal C_{L_{\mathcal N}}(\mathcal N)
         & \subseteq[-1,1]\cap\mathbb Z.
    \end{align}
    With the finite-supremum convention, $L_{\varnothing}=1$.
\end{lemma}

\begin{proof}\leanok
    \uses{lem:qca_blocked_neighborhood_subset_nearest_neighbor_sup}
    The finite supremum is strictly smaller than its successor, which is
    positive even when $\mathcal N$ is empty.
\end{proof}

\begin{lemma}[Nearest-neighbor propagation after a bounded blocking]
    \label{lem:qca_propagates_within_blocked_nearest_neighbor}
    \lean{SpinChain.PropagatesWithin.blocked_nearest_neighbor}
    \leanok
    \uses{def:qca_propagates_within, def:qca_blocked_automorphism}
    Let $d,L>0$.  Suppose that $\omega$ propagates within the finite
    neighborhood $\mathcal N$ and that $|n|<L$ for every $n\in\mathcal N$.
    For every finite blocked region $\Lambda$ and every observable
    $X\in\mathcal A^{(d^L)}$ supported in $\Lambda$, the observable
    $\omega^{[L]}(X)$ is supported in
    \begin{align}
        \Lambda+([-1,1]\cap\mathbb Z).
    \end{align}
\end{lemma}

\begin{proof}\leanok
    \uses{lem:qca_propagates_within_blocked, lem:qca_blocked_neighborhood_subset_nearest_neighbor, lem:qca_propagation_neighborhood_mono}
    Blocking first gives propagation within $\mathcal C_L(\mathcal N)$.
    The carry bound places this neighborhood inside
    $[-1,1]\cap\mathbb Z$, and enlargement of propagation neighborhoods gives
    the claim.
\end{proof}

\begin{lemma}[Finite propagation becomes nearest-neighbor]
    \label{thm:qca_finite_propagation_exists_blocked_nearest_neighbor}
    \lean{SpinChain.HasFinitePropagation.exists_blocked_nearest_neighbor}
    \leanok
    \uses{def:qca_finite_propagation, def:qca_propagates_within, def:qca_blocked_automorphism}
    Let $d>0$.  If $\omega$ has finite forward propagation, then there exists
    a positive integer $L$ such that, for every finite blocked region
    $\Lambda$ and every observable $X\in\mathcal A^{(d^L)}$ supported in
    $\Lambda$, the observable $\omega^{[L]}(X)$ is supported in
    \begin{align}
        \Lambda+([-1,1]\cap\mathbb Z).
    \end{align}
\end{lemma}

\begin{proof}\leanok
    \uses{lem:qca_propagates_within_blocked, lem:qca_blocked_neighborhood_canonical_nearest_neighbor, lem:qca_propagation_neighborhood_mono}
    For a finite propagation neighborhood $\mathcal N$, let
    $L=1+\sup_{n\in\mathcal N}|n|$.  This $L$ is positive, including when
    $\mathcal N$ is empty.  Blocking gives propagation within
    $\mathcal C_L(\mathcal N)$; the canonical carry inclusion and propagation
    monotonicity then give the nearest-neighbor bound.
\end{proof}

\begin{lemma}[A QCA becomes nearest-neighbor after blocking]
    \label{thm:qca_exists_blocked_nearest_neighbor}
    \lean{SpinChain.IsQCA.exists_blocked_nearest_neighbor}
    \leanok
    \uses{def:qca, def:qca_blocked_automorphism, def:qca_propagates_within}
    Let $d>0$ and let $\omega$ be a one-dimensional QCA on
    $\mathcal A^{(d)}$.  There exists a positive integer $L$ such that
    $\omega^{[L]}$ is a QCA on $\mathcal A^{(d^L)}$ and, for every finite
    blocked region $\Lambda$ and every observable
    $X\in\mathcal A^{(d^L)}$ supported in $\Lambda$, the observable
    $\omega^{[L]}(X)$ is supported in
    \begin{align}
        \Lambda+([-1,1]\cap\mathbb Z).
    \end{align}
\end{lemma}

\begin{proof}\leanok
    \uses{lem:qca_blocked, thm:qca_finite_propagation_exists_blocked_nearest_neighbor}
    Lemma~\ref{thm:qca_finite_propagation_exists_blocked_nearest_neighbor} applied
    to the finite propagation of $\omega$ at
    the canonical positive block length gives nearest-neighbor propagation; at the same length, site blocking
    preserves both finite propagation and translation covariance, hence the
    QCA property.
\end{proof}

\section{Support algebras of a quantum cellular automaton}

Lines~2292--2298 of \cite{Cirac2017MPU} specify the finite local
algebras, their canonical inclusions, translation covariance, and the locality
inclusion
\begin{align}
    \omega(\mathcal A_\Lambda)
     & \subseteq \mathcal A_{\Lambda+\mathcal N}.
\end{align}
Following \cite{Gross2012IndexTheory}, the finite-dimensional maps determined by this locality inclusion use the
adjacent-pair ordering convention.

\begin{definition}[Finite-region restriction]
    \label{def:qca_finite_region_restriction}
    \lean{SpinChain.PropagatesWithin.localRestriction}
    \leanok
    \uses{def:qca_local_algebra, def:qca_local_inclusion, def:qca_region_sumset, def:qca_quasi_local_observable, def:qca_propagates_within}
    Let $d>0$, let $\omega$ be a star-algebra automorphism of
    $\mathcal A^{(d)}$, and suppose that $\omega$ propagates within the finite
    neighborhood $\mathcal N$.  For every finite region $\Lambda$, define
    the finite-region restriction
    \begin{align}
        \omega_{\Lambda,\mathcal N}\colon\mathcal A_\Lambda
         & \longrightarrow\mathcal A_{\Lambda+\mathcal N}
    \end{align}
    to be the unique star-algebra homomorphism satisfying
    \begin{align}
        \iota_{\Lambda+\mathcal N}
        \bigl(\omega_{\Lambda,\mathcal N}(A)\bigr)
         & =\omega\bigl(\iota_\Lambda(A)\bigr)
        \qquad(A\in\mathcal A_\Lambda).
    \end{align}
\end{definition}

\begin{lemma}[Finite-region restriction is an embedding]
    \label{thm:qca_finite_region_restriction_embedding}
    \lean{SpinChain.PropagatesWithin.quasiLocalObservable_localRestriction,
        SpinChain.PropagatesWithin.localRestriction_injective}
    \leanok
    \uses{def:qca_finite_region_restriction}
    The finite-region restriction is injective, and for every
    $A\in\mathcal A_\Lambda$ one has
    \begin{align}
        \iota_{\Lambda+\mathcal N}
        \bigl(\omega_{\Lambda,\mathcal N}(A)\bigr)
         & =\omega\bigl(\iota_\Lambda(A)\bigr).
    \end{align}
\end{lemma}

\begin{proof}\leanok
    \uses{def:qca_finite_region_restriction}
    The defining equality follows from the locality inclusion and the
    injectivity of $\iota_{\Lambda+\mathcal N}$.  If two local observables
    have the same image under $\omega_{\Lambda,\mathcal N}$, their canonical
    images become equal after applying $\omega$.  Injectivity of $\omega$ and
    of $\iota_\Lambda$ then gives equality of the two local observables.
\end{proof}

\begin{lemma}[Enlargement of finite-region restrictions]
    \label{thm:qca_finite_region_restriction_enlargement}
    \lean{SpinChain.PropagatesWithin.localRestriction_localInclusion}
    \leanok
    \uses{def:qca_finite_region_restriction}
    If $\Lambda\subseteq\Gamma$, then for every
    $A\in\mathcal A_\Lambda$,
    \begin{align}
        \omega_{\Gamma,\mathcal N}
        \bigl(\iota_{\Lambda,\Gamma}(A)\bigr)
         & =\iota_{\Lambda+\mathcal N,\Gamma+\mathcal N}
        \bigl(\omega_{\Lambda,\mathcal N}(A)\bigr).
    \end{align}
\end{lemma}

\begin{proof}\leanok
    \uses{thm:qca_finite_region_restriction_embedding}
    Under the canonical inclusion into the quasi-local algebra,
    both expressions become $\omega(\iota_\Lambda(A))$; injectivity of the
    target inclusion gives the equality.
\end{proof}

\begin{lemma}[Translation of finite-region restrictions]
    \label{thm:qca_finite_region_restriction_translation}
    \lean{SpinChain.PropagatesWithin.localRestriction_translation}
    \leanok
    \uses{def:qca_translation_covariant, def:qca_finite_local_translation, def:qca_finite_region_restriction}
    Suppose in addition that $\omega$ is translation covariant.  For
    $a\in\mathbb Z$ and $A\in\mathcal A_\Lambda$,
    \begin{align}
        \iota_{(\Lambda+a)+\mathcal N}
        \bigl(\omega_{\Lambda+a,\mathcal N}
            (\tau_{a,\Lambda}(A))\bigr)
         & =\iota_{(\Lambda+\mathcal N)+a}
        \bigl(\tau_{a,\Lambda+\mathcal N}
            (\omega_{\Lambda,\mathcal N}(A))\bigr).
    \end{align}
    Here $(\Lambda+a)+\mathcal N=(\Lambda+\mathcal N)+a$.
\end{lemma}

\begin{proof}\leanok
    \uses{thm:qca_finite_region_restriction_embedding, lem:qca_quasi_local_translation_action}
    The left-hand side is
    $\omega(\tau_a(\iota_\Lambda(A)))$.  Translation covariance moves
    $\tau_a$ past $\omega$, and the finite-region translation formula gives
    the right-hand side.
\end{proof}

\begin{definition}[Bipartite coordinates for disjoint local algebras]
    \label{def:qca_bipartite_local_algebra_coordinates}
    \lean{SpinChain.Config.disjointUnionEquiv,
        SpinChain.bipartiteLocalAlgebraEquiv}
    \leanok
    \uses{def:qca_local_algebra}
    Let $\Gamma,\Delta\subset\mathbb Z$ be disjoint finite regions.  Writing
    $\mathcal C_\Gamma$ and $\mathcal C_\Delta$ for their configuration
    spaces, restriction to the two regions gives
    \begin{align}
        \mathcal C_{\Gamma\cup\Delta}
         & \simeq \mathcal C_\Gamma\times\mathcal C_\Delta,
    \end{align}
    with the $\Gamma$-configuration first.  Reindexing matrix rows and
    columns gives a star-algebra equivalence
    \begin{align}
        \beta_{\Gamma,\Delta}\colon
        \mathcal A_{\Gamma\cup\Delta}
         & \simeq
        M_{\mathcal C_\Gamma\times\mathcal C_\Delta}(\mathbb C).
    \end{align}
\end{definition}

\begin{lemma}[Canonical inclusions in bipartite coordinates]
    \label{thm:qca_bipartite_local_algebra_inclusions}
    \lean{SpinChain.bipartiteLocalAlgebraEquiv_localInclusion_left,
        SpinChain.bipartiteLocalAlgebraEquiv_localInclusion_right}
    \leanok
    \uses{def:qca_bipartite_local_algebra_coordinates, def:qca_local_inclusion}
    For $A\in\mathcal A_\Gamma$ and $B\in\mathcal A_\Delta$,
    \begin{align}
        \beta_{\Gamma,\Delta}(\iota_{\Gamma,\Gamma\cup\Delta}(A))
         & =A\otimes\mathbf 1,  \\
        \beta_{\Gamma,\Delta}(\iota_{\Delta,\Gamma\cup\Delta}(B))
         & =\mathbf 1\otimes B.
    \end{align}
\end{lemma}

\begin{proof}\leanok
    \uses{def:qca_bipartite_local_algebra_coordinates}
    Evaluating both matrices on pairs of configurations shows that, for the first local
    inclusion, restriction to $\Gamma$ gives the matrix coefficient of $A$,
    while equality of the complementary $\Delta$-configurations gives the
    matrix coefficient of the identity.  This is $A\otimes\mathbf 1$.
    Interchanging the two regions gives $\mathbf 1\otimes B$.
\end{proof}

\begin{definition}[Finite local image in bipartite coordinates]
    \label{def:qca_bipartite_finite_local_image}
    \lean{SpinChain.PropagatesWithin.localRestrictionRange,
        SpinChain.PropagatesWithin.localRestrictionRangeEquiv,
        SpinChain.StarSubalgebra.bipartiteReindex,
        SpinChain.PropagatesWithin.bipartiteLocalRestrictionRange}
    \leanok
    \uses{def:qca_finite_region_restriction, thm:qca_finite_region_restriction_embedding, def:qca_bipartite_local_algebra_coordinates}
    Define the finite local image by
    \begin{align}
        \mathcal R_{\Lambda,\mathcal N}
         & =\omega_{\Lambda,\mathcal N}(\mathcal A_\Lambda)
        \subseteq\mathcal A_{\Lambda+\mathcal N}.
    \end{align}
    It is a star-subalgebra star-isomorphic to $\mathcal A_\Lambda$.  If
    $\Lambda+\mathcal N=\Gamma\cup\Delta$ with
    $\Gamma\cap\Delta=\varnothing$, then
    \begin{align}
        \beta_{\Gamma,\Delta}
        (\mathcal R_{\Lambda,\mathcal N})
         & \subseteq
        M_{\mathcal C_\Gamma\times\mathcal C_\Delta}(\mathbb C)
    \end{align}
    is the bipartite matrix star-subalgebra to which the left and right
    support-algebra constructions apply.  No support algebra is defined in
    this statement.
\end{definition}

\begin{definition}[Nearest-neighbour pair regions]
    \label{def:qca_nearest_neighbor_pair_regions}
    \lean{SpinChain.evenPair,SpinChain.leftPair,SpinChain.rightPair}
    \leanok
    For $x\in\mathbb Z$, set
    \begin{align}
        E_x & =\{2x,2x+1\},   &
        L_x & =\{2x-1,2x\},   &
        P_x & =\{2x+1,2x+2\}.
        \label{eq:qca_pair_regions}
    \end{align}
    The order $L_x,P_x$ is the left--right factor order in equation
    \texttt{RR2x} of
    \cite[lines~1251--1266]{Gross2012IndexTheory}, and \cite[lines~2313--2318]{Cirac2017MPU} uses the same four-site
    pair
    localization.
\end{definition}

\begin{lemma}[Geometry of the nearest-neighbour pair image]
    \label{lem:qca_nearest_neighbor_pair_geometry}
    \lean{SpinChain.regionSumset_evenPair,
        SpinChain.disjoint_leftPair_rightPair,
        SpinChain.rightPair_eq_leftPair_add_one}
    \leanok
    \uses{def:qca_nearest_neighbor_pair_regions}
    For the neighbourhood $\mathcal N=\{-1,0,1\}$,
    \begin{align}
        E_x+\mathcal N & =L_x\cup P_x, &
        L_x\cap P_x    & =\varnothing, &
        P_x            & =L_{x+1}.
        \label{eq:qca_pair_geometry}
    \end{align}
\end{lemma}

\begin{proof}\leanok
    By (\ref{eq:qca_pair_regions}), adding $-1,0,1$ to the two
    points of $E_x$ gives
    \begin{align}
        E_x+\{-1,0,1\}
         & =\{2x-1,2x,2x+1,2x+2\}=L_x\cup P_x.
    \end{align}
    The largest point of $L_x$ is $2x$, while the smallest point of $P_x$
    is $2x+1$, so the two sets are disjoint.  Finally,
    $L_{x+1}=\{2x+1,2x+2\}=P_x$.
\end{proof}

\begin{definition}[Site-indexed support algebras]
    \label{def:qca_site_indexed_support_algebras}
    \lean{SpinChain.matrixToQuasiLocalObservable,
        SpinChain.PropagatesWithin.evenPairLocalImage,
        SpinChain.PropagatesWithin.evenSupportAlgebra,
        SpinChain.PropagatesWithin.oddSupportAlgebra,
        SpinChain.PropagatesWithin.embeddedEvenSupportAlgebra,
        SpinChain.PropagatesWithin.embeddedOddSupportAlgebra,
        SpinChain.PropagatesWithin.embeddedSupportAlgebra,
        SpinChain.PropagatesWithin.embeddedSupportAlgebra_even,
        SpinChain.PropagatesWithin.embeddedSupportAlgebra_odd}
    \leanok
    \uses{def:qca_bipartite_finite_local_image, def:qca_quasi_local_observable, def:qca_left_support_algebra, def:qca_right_support_algebra, def:qca_nearest_neighbor_pair_regions, lem:qca_nearest_neighbor_pair_geometry}
    Let $\omega$ be a star-algebra automorphism of the homogeneous
    quasi-local algebra which propagates within $\{-1,0,1\}$.  Write the
    finite image of $\mathcal A_{E_x}$ in the ordered $L_x\times P_x$
    coordinates as
    \begin{align}
        S_x
         & =\beta_{L_x,P_x}\bigl(\omega(\mathcal A_{E_x})\bigr)
        \subseteq
        M_{\mathcal C_{L_x}\times\mathcal C_{P_x}}(\mathbb C).
        \label{eq:qca_site_image}
    \end{align}
    Define
    \begin{align}
        \mathcal R_{2x}
         & =\operatorname{Spp}_{\mathrm L}(S_x)
        \subseteq M_{\mathcal C_{L_x}}(\mathbb C), \\
        \mathcal R_{2x+1}
         & =\operatorname{Spp}_{\mathrm R}(S_x)
        \subseteq M_{\mathcal C_{P_x}}(\mathbb C).
        \label{eq:qca_RR2x}
    \end{align}
    After the canonical finite-region embeddings, these form one family of
    star-subalgebras of the quasi-local algebra.  Here the matrix embedding is
    the composition
    \begin{align}
        M_{\mathcal C_\Lambda}(\mathbb C)
        \cong \mathcal A_\Lambda
        \longrightarrow \mathcal A.
    \end{align}
    Under this identification, the canonical local inclusion in
    \cite[lines~2292--2298]{Cirac2017MPU} is exactly the displayed
    composition.
    Formula (\ref{eq:qca_RR2x}) is equation \texttt{RR2x} in
    \cite[lines~1261--1266]{Gross2012IndexTheory}.  Cirac et al. use the
    homogeneous even factor in
    \cite[lines~2322--2328]{Cirac2017MPU}; the odd factor and unified family
    are given by the GNVW construction.

    While GNVW allow the cell size $d(y)$ to depend on $y$, as documented in \cite{gap:gnvw12_site_dependent_adjacent_generation_scope} the present
    definition treats a homogeneous chain with fixed cell size
    $d$.
\end{definition}

\begin{lemma}[Adjacent-pair support containment]
    \label{thm:qca_adjacent_pair_support_containment}
    \lean{SpinChain.PropagatesWithin.evenPairLocalImage_le_support_kroneckerSubmodule}
    \leanok
    \uses{def:qca_site_indexed_support_algebras, thm:qca_two_sided_support_containment}
    For every $x\in\mathbb Z$, the finite image satisfies
    \begin{align}
        S_x
         & \subseteq
        \mathcal R_{2x}\boxtimes\mathcal R_{2x+1}
        \subseteq
        M_{\mathcal C_{L_x}\times\mathcal C_{P_x}}(\mathbb C).
        \label{eq:qca_adjacent_pair_containment}
    \end{align}
    This is the homogeneous, full-matrix instance used in the two-sided
    containment at
    \cite[lines~1276--1278]{Gross2012IndexTheory} (see \cite{gap:gnvw12_site_dependent_adjacent_generation_scope} for the homogeneous-chain scope),
    whereas
    at \cite[lines~2322--2328]{Cirac2017MPU} Cirac et al. state only the weaker inclusion in the left support algebra
    tensored with the unrestricted right physical
    factor.
\end{lemma}

\begin{proof}\leanok
    \uses{thm:qca_two_sided_support_containment}
    By Lemma~\ref{thm:qca_two_sided_support_containment} applied to $S_x$ in
    the ordered coordinates of (\ref{eq:qca_site_image}), its left support
    is $\mathcal R_{2x}$ and its right support is
    $\mathcal R_{2x+1}$ by (\ref{eq:qca_RR2x}), which gives
    (\ref{eq:qca_adjacent_pair_containment}).
\end{proof}

\begin{lemma}[Physical support of the site-indexed support algebras]
    \label{thm:qca_site_indexed_support_regions}
    \lean{SpinChain.matrixToQuasiLocalObservable_supportedIn,
        SpinChain.PropagatesWithin.embeddedSupportAlgebra_even_supportedIn,
        SpinChain.PropagatesWithin.embeddedSupportAlgebra_odd_supportedIn}
    \leanok
    \uses{def:qca_site_indexed_support_algebras, def:qca_quasi_local_supported_in}
    If $\widehat{\mathcal R}_y$ denotes the canonical quasi-local copy of
    $\mathcal R_y$, then for every $x\in\mathbb Z$,
    \begin{align}
        \widehat{\mathcal R}_{2x}   & \subseteq\mathcal A_{L_x}, &
        \widehat{\mathcal R}_{2x+1} & \subseteq\mathcal A_{P_x}.
    \end{align}
    More generally, the canonical image of any matrix on a finite region
    $\Lambda$ is supported in $\Lambda$.
    This is the homogeneous specialization of the localization accompanying
    equation \texttt{RR2x} in \cite[lines~1261--1274]{Gross2012IndexTheory}, with \cite{gap:gnvw12_site_dependent_adjacent_generation_scope} recording the
    fixed-$d$
    scope.
\end{lemma}

\begin{proof}\leanok
    An element of $\widehat{\mathcal R}_{2x}$ is the canonical image of a
    matrix in $\mathcal R_{2x}\subseteq\mathcal A_{L_x}$ and is therefore
    supported in $L_x$.  The same argument places
    $\widehat{\mathcal R}_{2x+1}$ in $\mathcal A_{P_x}$.
\end{proof}

\begin{lemma}[Generation by two disjoint finite regions]
    \label{thm:qca_disjoint_region_generation}
    \lean{SpinChain.range_localObservable_mono,
        SpinChain.range_localObservable_union}
    \leanok
    \uses{def:qca_algebraic_local_algebra, def:qca_local_inclusion,
        def:qca_bipartite_local_algebra_coordinates,
        thm:qca_bipartite_local_algebra_inclusions,
        thm:qca_commuting_factor_kronecker_coordinates}
    If $\Lambda\subseteq\Gamma$ are finite regions, then
    $\iota_\Lambda(\mathcal A_\Lambda)
        \subseteq\iota_\Gamma(\mathcal A_\Gamma)$.  If the finite regions
    $\Lambda$ and $\Gamma$ are disjoint, then the two canonical images
    generate the image of their union,
    \begin{align}
        \iota_{\Lambda\cup\Gamma}(\mathcal A_{\Lambda\cup\Gamma})
         & =\iota_\Lambda(\mathcal A_\Lambda)
        \vee\iota_\Gamma(\mathcal A_\Gamma).
        \label{eq:qca_disjoint_region_generation}
    \end{align}
    This is the expansion of an observable on several cells into a product of
    observables on disjoint groups of cells used in \cite[lines~592--596]{Gross2012IndexTheory} and proved in \cite[lines~341--364]{SchumacherWerner2004QCA} for finite
    cell
    partitions.
\end{lemma}

\begin{proof}\leanok
    \uses{thm:qca_bipartite_local_algebra_inclusions,
        thm:qca_commuting_factor_kronecker_coordinates,
        thm:qca_mem_kronecker_submodule_iff}
    Monotonicity follows from
    $\iota_\Gamma\circ\iota_{\Lambda,\Gamma}=\iota_\Lambda$, and it gives the
    inclusion $\supseteq$ in
    (\ref{eq:qca_disjoint_region_generation}).  For the remaining inclusion,
    let $\mathcal V\subseteq\mathcal A_{\Lambda\cup\Gamma}$ be the
    star-subalgebra of those $A$ whose canonical image lies in the join on the
    right-hand side.  In the bipartite coordinates
    $\beta_{\Lambda,\Gamma}$, the canonical inclusion of
    $\mathcal A_\Lambda$ is $X\mapsto X\otimes\mathbf 1$ and that of
    $\mathcal A_\Gamma$ is $Y\mapsto\mathbf 1\otimes Y$, so
    $\beta_{\Lambda,\Gamma}(\mathcal V)$ contains
    $M_{\mathcal C_\Lambda}(\mathbb C)\otimes\mathbf 1$ and
    $\mathbf 1\otimes M_{\mathcal C_\Gamma}(\mathbb C)$.  The join of these two
    subalgebras has underlying vector space
    $M_{\mathcal C_\Lambda}(\mathbb C)\boxtimes
        M_{\mathcal C_\Gamma}(\mathbb C)$, and every matrix on
    $\mathcal C_\Lambda\times\mathcal C_\Gamma$ lies in that span.  Hence
    $\beta_{\Lambda,\Gamma}(\mathcal V)$ is everything, and
    $\beta_{\Lambda,\Gamma}$ is bijective, so
    $\mathcal V=\mathcal A_{\Lambda\cup\Gamma}$.
\end{proof}

\begin{lemma}[Even pairs cover the chain]
    \label{lem:qca_even_pair_cover}
    \lean{SpinChain.disjoint_evenPair_of_ne,
        SpinChain.subset_biUnion_evenPair}
    \leanok
    \uses{def:qca_nearest_neighbor_pair_regions}
    For $x\ne y$ the pairs $E_x$ and $E_y$ are disjoint, and every finite
    region $\Lambda$ satisfies
    \begin{align}
        \Lambda & \subseteq\bigcup_{x\in X_\Lambda}E_x,
        \label{eq:qca_even_pair_cover}
    \end{align}
    where $X_\Lambda=\{\lfloor i/2\rfloor\mid i\in\Lambda\}$.
\end{lemma}

\begin{proof}\leanok
    \uses{def:qca_nearest_neighbor_pair_regions}
    Both statements are integer arithmetic: the four sites
    $2x,2x+1,2y,2y+1$ are distinct for $x\ne y$, and every $i$ satisfies
    $i=2\lfloor i/2\rfloor$ or $i=2\lfloor i/2\rfloor+1$.
\end{proof}

\begin{lemma}[Generation of the algebraic local algebra by the pair algebras]
    \label{thm:qca_even_pair_algebraic_generation}
    \lean{SpinChain.range_localObservable_biUnion_evenPair_le,
        SpinChain.iSup_range_localObservable_evenPair}
    \leanok
    \uses{def:qca_nearest_neighbor_pair_regions,
        def:qca_algebraic_local_algebra,
        thm:qca_disjoint_region_generation, lem:qca_even_pair_cover,
        thm:qca_local_observable_finite_support}
    The pair algebras generate the algebraic local algebra,
    \begin{align}
        \bigvee_{x\in\mathbb Z}\iota_{E_x}(\mathcal A_{E_x})
         & =\mathcal A_{\mathrm{loc}}.
        \label{eq:qca_even_pair_algebraic_generation}
    \end{align}
    The completion is not involved: this is an equality of algebras of local
    observables.
\end{lemma}

\begin{proof}\leanok
    \uses{thm:qca_disjoint_region_generation, lem:qca_even_pair_cover,
        thm:qca_local_observable_finite_support}
    For a finite set $X$ of pair labels, induction on $X$ gives
    \begin{align}
        \iota_{\bigcup_{x\in X}E_x}
        \Bigl(\mathcal A_{\bigcup_{x\in X}E_x}\Bigr)
         & \subseteq\bigvee_{x\in\mathbb Z}
        \iota_{E_x}(\mathcal A_{E_x}).
        \label{eq:qca_even_pair_generation_induction}
    \end{align}
    For $X=\varnothing$ the region is empty and contained in $E_0$, so
    monotonicity applies.  For $X\cup\{x_0\}$ with $x_0\notin X$, the pair
    $E_{x_0}$ is disjoint from $\bigcup_{x\in X}E_x$, so
    (\ref{eq:qca_disjoint_region_generation}) splits the left-hand side into
    the join of $\iota_{E_{x_0}}(\mathcal A_{E_{x_0}})$ with the algebra
    handled by the induction hypothesis.  Every local observable is supported
    in some finite region $\Lambda$, which is covered by
    (\ref{eq:qca_even_pair_cover}); monotonicity and
    (\ref{eq:qca_even_pair_generation_induction}) place it in the join.
\end{proof}

\begin{definition}[Join of the site-indexed support algebras]
    \label{def:qca_support_algebra_join}
    \lean{SpinChain.PropagatesWithin.supportAlgebraJoin,
        SpinChain.PropagatesWithin.embeddedEvenSupportAlgebra_le_supportAlgebraJoin,
        SpinChain.PropagatesWithin.embeddedOddSupportAlgebra_le_supportAlgebraJoin}
    \leanok
    \uses{def:qca_site_indexed_support_algebras}
    For a star-algebra automorphism $\omega$ of the quasi-local algebra
    propagating within $\{-1,0,1\}$, define
    \begin{align}
        G & =\bigvee_{y\in\mathbb Z}\widehat{\mathcal R}_y
        \subseteq\mathcal A.
        \label{eq:qca_support_algebra_join}
    \end{align}
    Each $\widehat{\mathcal R}_y$ is contained in $G$.  This is the algebraic
    join associated with the algebra generated by all $\mathcal R_y$ in
    \cite[lines~1276--1279]{Gross2012IndexTheory}.  No claim that $G$ is
    proper or that it fails to be norm closed is made here.
\end{definition}

\begin{lemma}[Images of the pair algebras lie in the join]
    \label{thm:qca_evenPair_image_in_support_algebra_join}
    \lean{SpinChain.PropagatesWithin.quasiLocalObservable_evenPair_mem_supportAlgebraJoin,
        SpinChain.PropagatesWithin.algebraicToQuasiLocal_mem_supportAlgebraJoin}
    \leanok
    \uses{def:qca_support_algebra_join,
        def:qca_bipartite_finite_local_image,
        thm:qca_adjacent_pair_support_containment,
        thm:qca_even_pair_algebraic_generation}
    For every $x\in\mathbb Z$,
    \begin{align}
        \omega\bigl(j_{E_x}(\mathcal A_{E_x})\bigr) & \subseteq G,
        \label{eq:qca_evenPair_image_in_support_algebra_join}
    \end{align}
    and consequently
    $\omega\bigl(j(\mathcal A_{\mathrm{loc}})\bigr)\subseteq G$.
    Inclusion (\ref{eq:qca_evenPair_image_in_support_algebra_join}) is the
    containment
    $\alpha(\mathcal A_{2x}\otimes\mathcal A_{2x+1})\subseteq
        \mathcal R_{2x}\otimes\mathcal R_{2x+1}$
    at \cite[lines~1276--1278]{Gross2012IndexTheory}, specialized as in \cite{gap:gnvw12_site_dependent_adjacent_generation_scope} to the homogeneous
    fixed-$d$
    setting.
\end{lemma}

\begin{proof}\leanok
    \uses{thm:qca_adjacent_pair_support_containment,
        thm:qca_commuting_factor_kronecker_coordinates,
        thm:qca_bipartite_local_algebra_inclusions,
        thm:qca_even_pair_algebraic_generation}
    Let $S_x$ be the finite image of $\mathcal A_{E_x}$ in the ordered
    $L_x\times P_x$ coordinates.  By
    (\ref{eq:qca_adjacent_pair_containment}), $S_x$ is contained in
    $\mathcal R_{2x}\boxtimes\mathcal R_{2x+1}$, which is the underlying vector
    space of the join of $\mathcal R_{2x}\otimes\mathbf 1$ and
    $\mathbf 1\otimes\mathcal R_{2x+1}$.  The canonical quasi-local map on
    $L_x\cup P_x$ sends $X\otimes\mathbf 1$ to the canonical embedding of
    $X\in\mathcal R_{2x}$ and $\mathbf 1\otimes Y$ to the canonical embedding
    of $Y\in\mathcal R_{2x+1}$, so it maps that join into $G$.  It sends the
    image of $A\in\mathcal A_{E_x}$ to $\omega(j_{E_x}(A))$, which proves
    (\ref{eq:qca_evenPair_image_in_support_algebra_join}).  By
    (\ref{eq:qca_even_pair_algebraic_generation}), every local observable lies
    in the algebra generated by the $\iota_{E_x}(\mathcal A_{E_x})$, and
    $\omega\circ j$ is a star-algebra homomorphism, so its image lies in $G$.
\end{proof}

\begin{lemma}[Norm-closed generation by the support algebras]
    \label{thm:qca_support_algebra_topological_generation}
    \lean{SpinChain.PropagatesWithin.supportAlgebraJoin_topologicalClosure_eq_top}
    \leanok
    \uses{def:qca_support_algebra_join,
        thm:qca_evenPair_image_in_support_algebra_join,
        thm:qca_quasi_local_density}
    The site-indexed support algebras generate the quasi-local algebra in the
    norm-closed sense,
    \begin{align}
        \overline{G}^{\lVert\cdot\rVert} & =\mathcal A.
        \label{eq:qca_support_algebra_topological_generation}
    \end{align}
    In \cite[lines~1276--1279]{Gross2012IndexTheory} the algebras
    $\mathcal R_y$ are said to generate an algebra containing
    $\alpha(\mathcal A(\mathbb Z))$, which equals $\mathcal A(\mathbb Z)$
    because $\alpha$ is an automorphism.  Here this generation statement is
    expressed by the norm-closure equality
    (\ref{eq:qca_support_algebra_topological_generation}), specialized as in \cite{gap:gnvw12_site_dependent_adjacent_generation_scope} to the homogeneous
    fixed-$d$
    setting.
\end{lemma}

\begin{proof}\leanok
    \uses{thm:qca_evenPair_image_in_support_algebra_join,
        thm:qca_quasi_local_density}
    An injective star-algebra homomorphism of complex C-star algebras is
    isometric, so $\omega$ is continuous and
    $\omega^{-1}(\overline G)$ is closed.  By
    (\ref{eq:qca_evenPair_image_in_support_algebra_join}) it contains
    $j(\mathcal A_{\mathrm{loc}})$, which is dense, so
    $\omega^{-1}(\overline G)=\mathcal A$.  Since $\omega$ is surjective,
    every $z\in\mathcal A$ equals $\omega(\omega^{-1}(z))\in\overline G$.
\end{proof}

\begin{corollary}[Commutant of the site-indexed support algebras]
    \label{cor:qca_support_algebra_commutant}
    \lean{SpinChain.PropagatesWithin.commute_of_forall_commute_embeddedSupportAlgebra}
    \leanok
    \uses{def:qca_support_algebra_join,
        thm:qca_support_algebra_topological_generation}
    If $z\in\mathcal A$ satisfies $[z,R]=0$ for every $y\in\mathbb Z$ and
    every $R\in\widehat{\mathcal R}_y$, then $[z,A]=0$ for every
    $A\in\mathcal A$.  This is the step from commutation with all
    $\mathcal R_y$ to commutation in the quasi-local algebra at \cite[lines~1279--1282]{Gross2012IndexTheory}, under the homogeneous scope recorded in \cite{gap:gnvw12_site_dependent_adjacent_generation_scope} for fixed
    cell
    size.
\end{corollary}

\begin{proof}\leanok
    \uses{def:qca_support_algebra_join,
        thm:qca_support_algebra_topological_generation}
    The elements commuting with $z$ and with $z^*$ form a norm-closed
    star-subalgebra.  Each $\widehat{\mathcal R}_y$ is star closed, so the
    hypothesis applied to $R^*$ gives $[z^*,R]=0$ as well; hence that
    star-subalgebra contains every $\widehat{\mathcal R}_y$, therefore
    contains $G$, and being closed it contains
    $\overline G=\mathcal A$ by
    (\ref{eq:qca_support_algebra_topological_generation}).
\end{proof}

\begin{definition}[Tripartite finite-region coordinates]
    \label{def:qca_tripartite_local_coordinates}
    \lean{SpinChain.Config.disjointTripleEquiv,
        SpinChain.tripartiteLocalAlgebraEquiv}
    \leanok
    \uses{def:qca_bipartite_local_algebra_coordinates}
    For pairwise disjoint finite regions $\Lambda,\Gamma,\Delta$, successive
    restriction in left--middle--right order gives
    \begin{align}
        \mathcal C_{(\Lambda\cup\Gamma)\cup\Delta}
         & \cong(\mathcal C_\Lambda\times\mathcal C_\Gamma)
        \times\mathcal C_\Delta,                                \\
        \mathcal A_{(\Lambda\cup\Gamma)\cup\Delta}
         & \cong M_{(\mathcal C_\Lambda\times\mathcal C_\Gamma)
                \times\mathcal C_\Delta}(\mathbb C).
    \end{align}
\end{definition}

\begin{lemma}[Geometry of consecutive pair images]
    \label{lem:qca_consecutive_pair_geometry}
    \lean{SpinChain.disjoint_leftPair_union_rightPair_rightPair_add_one,
        SpinChain.leftPair_union_rightPair_union_rightPair_add_one}
    \leanok
    \uses{def:qca_nearest_neighbor_pair_regions}
    Let $T_x=(L_x\cup P_x)\cup P_{x+1}$. The three regions
    $L_x,P_x,P_{x+1}$ are pairwise disjoint, and
    \begin{align}
        T_x & =\{2x-1,2x,2x+1,2x+2,2x+3,2x+4\}.
    \end{align}
\end{lemma}

\begin{proof}\leanok
    \uses{def:qca_nearest_neighbor_pair_regions, lem:qca_nearest_neighbor_pair_geometry}
    Unfolding the three pair definitions and comparing their six sites gives both disjointness relations.
\end{proof}

\begin{lemma}[Canonical inclusions in tripartite coordinates]
    \label{thm:qca_consecutive_pair_tripartite_coordinates}
    \lean{SpinChain.tripartiteLocalAlgebraEquiv_localInclusion_left,
        SpinChain.tripartiteLocalAlgebraEquiv_localInclusion_right}
    \leanok
    \uses{def:qca_tripartite_local_coordinates,
        def:qca_local_inclusion}
    In the coordinates of Definition~\ref{def:qca_tripartite_local_coordinates},
    the canonical inclusion from $\Lambda\cup\Gamma$ is
    $X\mapsto X\otimes\mathbf 1$, and the canonical inclusion from
    $\Gamma\cup\Delta$ is $Y\mapsto\mathbf 1\otimes Y$, with the ordered
    factors $\Lambda,\Gamma,\Delta$.
\end{lemma}

\begin{proof}\leanok
    \uses{def:qca_tripartite_local_coordinates,
        def:qca_bipartite_local_algebra_coordinates,
        def:qca_local_inclusion}
    Restricting each reconstructed configuration to the retained region and its
    complement shows that in the first inclusion the complementary coordinates are
    equal precisely on $\Delta$; in the second they are equal precisely on
    $\Lambda$.  The two matrix entries are therefore the left and right
    overlapping lifts, respectively.
\end{proof}

\begin{lemma}[Commutation across the common physical pair]
    \label{thm:qca_consecutive_support_algebras_commute}
    \lean{
        SpinChain.PropagatesWithin.embeddedOddSupport_commute_embeddedEvenSupport_add_one}
    \leanok
    \uses{def:qca_site_indexed_support_algebras}
    For every $x\in\mathbb Z$, the canonically embedded support algebras
    $\widehat{\mathcal R}_{2x+1}$ and
    $\widehat{\mathcal R}_{2x+2}$ commute elementwise.
\end{lemma}

\begin{proof}\leanok
    \uses{lem:qca_consecutive_pair_geometry,
        thm:qca_consecutive_pair_tripartite_coordinates,
        thm:qca_finite_region_restriction_embedding,
        thm:qca_disjoint_quasi_local_commute,
        thm:qca_overlapping_support_algebras_commute}
    The source regions $E_x$ and $E_{x+1}$ are disjoint, so their canonical
    quasi-local observables commute.  Multiplicativity of $\omega$ preserves
    this equality.  Embed the two finite images into $\mathcal A_{T_x}$ and
    use the tripartite coordinates above to obtain
    \begin{align}
        (X\otimes\mathbf 1)(\mathbf 1\otimes Y)
         & =(\mathbf 1\otimes Y)(X\otimes\mathbf 1).
    \end{align}
    The overlapping-support theorem, in the order
    $\operatorname{Spp}_{\mathrm R}(S_x)$ followed by
    $\operatorname{Spp}_{\mathrm L}(S_{x+1})$, gives the claim.
\end{proof}

\begin{definition}[Physical region of a site-indexed support algebra]
    \label{def:qca_support_algebra_region}
    \lean{SpinChain.supportRegion,
        SpinChain.supportRegion_even,
        SpinChain.supportRegion_odd}
    \leanok
    \uses{def:qca_nearest_neighbor_pair_regions}
    Define $Q_{2x}=L_x$ and $Q_{2x+1}=P_x$.
\end{definition}

\begin{lemma}[Unified physical localization]
    \label{thm:qca_unified_support_algebra_region}
    \lean{SpinChain.PropagatesWithin.embeddedSupportAlgebra_supportedIn}
    \leanok
    \uses{def:qca_support_algebra_region,
        def:qca_site_indexed_support_algebras}
    Every element of $\widehat{\mathcal R}_y$ is supported in $Q_y$.
\end{lemma}

\begin{proof}\leanok
    \uses{def:qca_support_algebra_region,
        thm:qca_site_indexed_support_regions}
    Splitting according to the parity of $y$ and applying the corresponding
    localization of $\widehat{\mathcal R}_{2x}$ or
    $\widehat{\mathcal R}_{2x+1}$ gives the three-site support.
\end{proof}

\begin{lemma}[Classification of overlapping support regions]
    \label{thm:qca_support_region_overlap_classification}
    \lean{SpinChain.PropagatesWithin.supportRegion_disjoint_or_consecutive}
    \leanok
    \uses{def:qca_support_algebra_region}
    If $y\ne z$, then $Q_y$ and $Q_z$ are disjoint unless
    \begin{align}
        (y,z) & =(2x+1,2x+2)
    \end{align}
    for some $x$, or unless this ordered pair is reversed.
\end{lemma}

\begin{proof}\leanok
    \uses{def:qca_support_algebra_region}
    Writing each index in its even or odd form shows that two regions of equal parity
    are disjoint.  In the mixed-parity cases, the defining two-site sets
    intersect exactly for the displayed consecutive odd--even pair.
\end{proof}

\begin{lemma}[Pairwise commutation of distinct support algebras]
    \label{thm:qca_site_indexed_support_algebras_commute}
    \lean{SpinChain.PropagatesWithin.embeddedSupportAlgebra_pairwise_commute}
    \leanok
    \uses{def:qca_site_indexed_support_algebras}
    Distinct members of the canonically embedded family commute elementwise:
    \begin{align}
        [X,Y] & =0,
              & X\in\widehat{\mathcal R}_y,\quad
        Y\in\widehat{\mathcal R}_z,\quad y\ne z.
    \end{align}
    This is the homogeneous fixed-$d$ specialization of the
    consecutive/disjoint argument in
    \cite[lines~1270--1274]{Gross2012IndexTheory}, using Lemma
    \texttt{sppcomm} from lines~1221--1246; compare
    \cite[lines~1174--1194]{SchumacherWerner2004QCA}.  The four-site localization is established in
    \cite[lines~2313--2329]{Cirac2017MPU}; see
    \cite{gap:gnvw12_site_dependent_adjacent_generation_scope} for the
    homogeneous scope of the adjacent-generation theorem.
\end{lemma}

\begin{proof}\leanok
    \uses{thm:qca_support_region_overlap_classification,
        thm:qca_unified_support_algebra_region,
        thm:qca_consecutive_support_algebras_commute,
        thm:qca_disjoint_quasi_local_commute}
    In the exceptional odd--even order apply
    Lemma~\ref{thm:qca_consecutive_support_algebras_commute}; the reversed
    order follows by symmetry.  Every remaining ordered pair has disjoint
    physical supports, so disjoint quasi-locality applies.
\end{proof}

\begin{lemma}[Two-site translation of the pair regions]
    \label{lem:qca_pair_region_two_site_translation}
    \lean{SpinChain.translateRegion_evenPair,
        SpinChain.translateRegion_leftPair,
        SpinChain.translateRegion_rightPair}
    \leanok
    \uses{def:qca_nearest_neighbor_pair_regions}
    Translation by two sites advances each of the three two-site regions by
    one index:
    \begin{align}
        E_x+2 & =E_{x+1}, &
        L_x+2 & =L_{x+1}, &
        P_x+2 & =P_{x+1}.
        \label{eq:qca_pair_region_two_site_translation}
    \end{align}
\end{lemma}

\begin{proof}\leanok
    \uses{def:qca_nearest_neighbor_pair_regions}
    Add $2$ to each of the two points of (\ref{eq:qca_pair_regions}) and
    compare with the same set at the index $x+1$.
\end{proof}

\begin{lemma}[Matrix coefficients as compressions by matrix units]
    \label{lem:qca_support_coefficient_compression}
    \lean{Matrix.leftKroneckerEmbed_bipartiteSlice_eq_sum,
        Matrix.rightKroneckerEmbed_operatorBlock_eq_sum}
    \leanok
    \uses{def:qca_left_coefficient_set, def:qca_right_coefficient_set}
    Let $X\in M_{I\times J}(\mathbb C)$ and write $e_{rs}$ for the matrix
    units of either factor.  Then
    \begin{align}
        X^{\mathrm L}_{rs}\otimes\mathbf 1
         & =\sum_{t\in J}
        (\mathbf 1\otimes e_{tr})\,X\,(\mathbf 1\otimes e_{st}),
        \label{eq:qca_left_coefficient_compression} \\
        \mathbf 1\otimes X^{\mathrm R}_{ij}
         & =\sum_{k\in I}
        (e_{ki}\otimes\mathbf 1)\,X\,(e_{jk}\otimes\mathbf 1).
        \label{eq:qca_right_coefficient_compression}
    \end{align}
    Every coefficient of the expansion used at
    \cite[lines~1188--1191]{Gross2012IndexTheory} and
    \cite[lines~1137--1169]{SchumacherWerner2004QCA} is therefore a sum of
    products formed inside the ambient algebra.
\end{lemma}

\begin{proof}\leanok
    Entrywise comparison proves both formulas: on the left of
    (\ref{eq:qca_left_coefficient_compression}), the entry at
    $((i,t),(j,u))$ equals $X((i,r),(j,s))$ when $t=u$ and vanishes
    otherwise.  On the right, the summand indexed by $t$ contributes
    $X((i,r),(j,s))$ exactly when both second coordinates equal $t$, so the
    sum over $t$ reproduces the same entry.  Exchanging the roles of the two factors yields~(\ref{eq:qca_right_coefficient_compression}).
\end{proof}

\begin{lemma}[Quasi-local form of the coefficient compressions]
    \label{lem:qca_quasi_local_coefficient_compression}
    \lean{SpinChain.matrixToQuasiLocalObservable_bipartiteSlice_eq_sum,
        SpinChain.matrixToQuasiLocalObservable_operatorBlock_eq_sum}
    \leanok
    \uses{lem:qca_support_coefficient_compression,
        def:qca_bipartite_local_algebra_coordinates,
        def:qca_quasi_local_observable}
    Let $\Gamma$ and $\Delta$ be disjoint finite regions and let
    $A\in\mathcal A_{\Gamma\cup\Delta}$, written in the ordered
    $\Gamma\times\Delta$ coordinates.  For configurations
    $r,s\in\mathcal C_\Delta$ and $i,j\in\mathcal C_\Gamma$, the canonical
    quasi-local copies of the two coefficient families are
    \begin{align}
        \widehat{A^{\mathrm L}_{rs}}
         & =\sum_{t\in\mathcal C_\Delta}
        \widehat{e^\Delta_{tr}}\ \widehat A\ \widehat{e^\Delta_{st}},
        \label{eq:qca_quasi_local_left_compression} \\
        \widehat{A^{\mathrm R}_{ij}}
         & =\sum_{k\in\mathcal C_\Gamma}
        \widehat{e^\Gamma_{ki}}\ \widehat A\ \widehat{e^\Gamma_{jk}},
        \label{eq:qca_quasi_local_right_compression}
    \end{align}
    where a hat denotes the canonical image in the quasi-local algebra.  Both
    right-hand sides are formed inside the quasi-local algebra alone.
\end{lemma}

\begin{proof}\leanok
    \uses{lem:qca_support_coefficient_compression,
        thm:qca_bipartite_local_algebra_inclusions,
        lem:qca_quasi_local_observable}
    In the bipartite coordinates the inclusion of $\mathcal A_\Gamma$ is
    $X\mapsto X\otimes\mathbf 1$ and the inclusion of $\mathcal A_\Delta$ is
    $Y\mapsto\mathbf 1\otimes Y$.  Passing
    (\ref{eq:qca_left_coefficient_compression}) and
    (\ref{eq:qca_right_coefficient_compression}) through the inverse of the
    bipartite coordinate equivalence, then apply the canonical map of
    $\mathcal A_{\Gamma\cup\Delta}$ into the quasi-local algebra, under which
    a finite-region inclusion leaves the quasi-local observable unchanged.
\end{proof}

\begin{lemma}[Two-site translation covariance of the support algebras]
    \label{thm:qca_site_indexed_support_algebras_two_site_covariance}
    \lean{SpinChain.PropagatesWithin.embeddedEvenSupportAlgebra_map_quasiLocalTranslation_two,
        SpinChain.PropagatesWithin.embeddedOddSupportAlgebra_map_quasiLocalTranslation_two,
        SpinChain.PropagatesWithin.embeddedSupportAlgebra_map_quasiLocalTranslation_two,
        SpinChain.PropagatesWithin.embeddedSupportAlgebra_map_quasiLocalTranslation_two_even,
        SpinChain.PropagatesWithin.embeddedSupportAlgebra_map_quasiLocalTranslation_two_odd}
    \leanok
    \uses{def:qca_site_indexed_support_algebras, def:qca_translation_covariant,
        def:qca_quasi_local_translation}
    Let $\omega$ be a star-algebra automorphism of the homogeneous quasi-local
    algebra which propagates within $\{-1,0,1\}$ and commutes with every
    lattice translation.  Then, for every $y\in\mathbb Z$,
    \begin{align}
        \tau_2\bigl(\widehat{\mathcal R}_y\bigr)
         & =\widehat{\mathcal R}_{y+2},
        \label{eq:qca_two_site_covariance}
    \end{align}
    and in the two parities
    \begin{align}
        \tau_2\bigl(\widehat{\mathcal R}_{2x}\bigr)
         & =\widehat{\mathcal R}_{2x+2}, &
        \tau_2\bigl(\widehat{\mathcal R}_{2x+1}\bigr)
         & =\widehat{\mathcal R}_{2x+3}.
        \label{eq:qca_two_site_covariance_parities}
    \end{align}
    The displacement is two sites of the original chain (one blocked cell); the
    shift automorphism in \cite[lines~1267--1268]{Gross2012IndexTheory} shows
    that single-site translation exchanges the inequivalent even and odd support
    algebras.

    Covariance of the support algebras under translation is the identity
    displayed at
    \cite[lines~1218--1226]{SchumacherWerner2004QCA} for the algebras defined
    at \cite[lines~1199--1206]{SchumacherWerner2004QCA}, and the ordered even
    and odd factors are equation \texttt{RR2x} at
    \cite[lines~1251--1266]{Gross2012IndexTheory}, while \cite[lines~2292--2298]{Cirac2017MPU} sets up the translation-covariant
    fixed-$d$
    chain.

    As recorded in \cite{gap:sw04_support_covariance_scope} and \cite{gap:gnvw12_site_dependent_adjacent_generation_scope}, the statement treats the one-dimensional translation-invariant
    setting with homogeneous cell size
    $d$.
\end{lemma}

\begin{proof}\leanok
    \uses{lem:qca_quasi_local_coefficient_compression,
        lem:qca_pair_region_two_site_translation,
        def:qca_site_indexed_support_algebras,
        def:qca_quasi_local_translation,
        thm:qca_finite_region_restriction_embedding}
    For a given $x$, let $\Lambda=E_x$, $\Gamma=L_x$, and $\Delta=P_x$.  A
    finite-region observable belongs to the finite local image exactly when
    its quasi-local copy lies in the image of $\mathcal A_\Lambda$ under
    $\omega$.  Hence, by
    (\ref{eq:qca_quasi_local_left_compression}), the canonical quasi-local
    copy of $\mathcal R_{2x}$ is generated by the elements
    \begin{align}
        \sum_{t\in\mathcal C_\Delta}
        \widehat{e^\Delta_{tr}}\ \omega\bigl(\widehat B\bigr)\
        \widehat{e^\Delta_{st}},
        \qquad B\in\mathcal A_\Lambda,\quad r,s\in\mathcal C_\Delta,
        \label{eq:qca_left_support_generators}
    \end{align}
    and the copy of $\mathcal R_{2x+1}$ is generated by the analogous
    elements built from (\ref{eq:qca_quasi_local_right_compression}) with
    $\Gamma$ in place of $\Delta$.  Since a star-algebra automorphism carries
    the algebra generated by a set to the algebra generated by the image of
    that set, it suffices to translate the two generating families.

    Translation carries a matrix unit of a region to the matrix unit of the
    translated region at the translated configurations, and covariance gives
    $\tau_2(\omega(\widehat B))=\omega(\tau_2(\widehat B))$, whose right-hand
    side is $\omega$ applied to the quasi-local copy of the translated
    observable.  Reindexing the summation over $\mathcal C_\Delta$ along the
    configuration translation therefore turns
    (\ref{eq:qca_left_support_generators}) into the same family for the
    regions $\Lambda+2$ and $\Delta+2$, and the translation is a bijection
    between the two families.  By
    (\ref{eq:qca_pair_region_two_site_translation}) those regions are
    $E_{x+1}$ and $P_{x+1}$, so the translated family generates
    $\widehat{\mathcal R}_{2x+2}$.  The odd case is identical with $\Gamma$
    in place of $\Delta$ and $L_{x+1}$ in place of $P_{x+1}$.

    Finally, splitting $y$ by parity and using $2x+2=2(x+1)$ and
    $2x+3=2(x+1)+1$ gives (\ref{eq:qca_two_site_covariance}) and
    (\ref{eq:qca_two_site_covariance_parities}).
\end{proof}

\section{Quantum cellular automata from finite-local data}

In \cite[lines~2300--2306]{Cirac2017MPU}, eventual constancy of finite-chain
conjugation defines a norm-preserving star-algebra automorphism of
$\mathcal A_{\mathrm{loc}}$ satisfying the support bound
$\mathcal N=[-4D^4,4D^4]\cap\mathbb Z$ and translation covariance, extending
uniquely to $\mathcal A$. The following statements develop this direct-limit
completion from compatible forward and inverse finite-local actions.

\begin{definition}[Compatible finite-local automorphism data]
    \label{def:qca_compatible_local_automorphism}
    \lean{SpinChain.CompatibleLocalAutomorphism}
    \leanok
    \uses{thm:qca_algebraic_local_algebra_universal}
    For each finite region $\Lambda$, let
    \begin{align}
        f_\Lambda,g_\Lambda\colon\mathcal A_\Lambda
         & \longrightarrow\mathcal A_{\mathrm{loc}}
    \end{align}
    be star-algebra homomorphisms such that, whenever
    $\Lambda\subseteq\Gamma$,
    \begin{align}
        f_\Gamma\circ\iota_{\Lambda,\Gamma} & =f_\Lambda, &
        g_\Gamma\circ\iota_{\Lambda,\Gamma} & =g_\Lambda.
    \end{align}
    Let $f,g\colon\mathcal A_{\mathrm{loc}}\to
        \mathcal A_{\mathrm{loc}}$ be the maps induced by these compatible
    families.  Compatible finite-local automorphism data require
    \begin{align}
        g(f_\Lambda(A)) & =\iota_\Lambda(A), &
        f(g_\Lambda(A)) & =\iota_\Lambda(A)
    \end{align}
    for every $A\in\mathcal A_\Lambda$.
\end{definition}

\begin{definition}[Induced algebraic-local homomorphisms]
    \label{def:qca_compatible_local_homs}
    \lean{SpinChain.CompatibleLocalAutomorphism.forwardHom,
        SpinChain.CompatibleLocalAutomorphism.inverseHom}
    \leanok
    \uses{def:qca_compatible_local_automorphism, thm:qca_algebraic_local_algebra_universal}
    The compatible families induce star-algebra homomorphisms
    \begin{align}
        f,g\colon\mathcal A_{\mathrm{loc}}
         & \longrightarrow\mathcal A_{\mathrm{loc}}.
    \end{align}
\end{definition}

\begin{proof}\leanok
    \uses{thm:qca_algebraic_local_algebra_universal}
    The star-algebra universal property applied to each compatible family produces both homomorphisms.
\end{proof}

\begin{lemma}[Evaluation of the induced homomorphisms]
    \label{lem:qca_compatible_local_homs_evaluation}
    \lean{SpinChain.CompatibleLocalAutomorphism.forwardHom_localObservable,
        SpinChain.CompatibleLocalAutomorphism.inverseHom_localObservable}
    \leanok
    \uses{def:qca_compatible_local_homs}
    For every finite region $\Lambda$ and $A\in\mathcal A_\Lambda$,
    \begin{align}
        f(\iota_\Lambda(A)) & =f_\Lambda(A), &
        g(\iota_\Lambda(A)) & =g_\Lambda(A).
    \end{align}
\end{lemma}

\begin{proof}\leanok
    \uses{thm:qca_algebraic_local_algebra_universal}
    These are the evaluation identities given by the universal property.
\end{proof}

\begin{definition}[Induced algebraic-local equivalence]
    \label{def:qca_compatible_algebraic_equiv}
    \lean{SpinChain.CompatibleLocalAutomorphism.algebraicEquiv}
    \leanok
    \uses{def:qca_compatible_local_automorphism, def:qca_compatible_local_homs}
    Compatible finite-local automorphism data determine a star-algebra
    equivalence
    \begin{align}
        F\colon\mathcal A_{\mathrm{loc}}
         & \simeq\mathcal A_{\mathrm{loc}}
    \end{align}
    whose forward and inverse maps are $f$ and $g$, respectively.
\end{definition}

\begin{proof}\leanok
    \uses{def:qca_compatible_local_automorphism, lem:qca_compatible_local_homs_evaluation}
    By induction on finite-region representatives in the algebraic direct limit,
    the two inverse assumptions give
    \begin{align}
        g(f(\iota_\Lambda(A))) & =g(f_\Lambda(A))=\iota_\Lambda(A), \\
        f(g(\iota_\Lambda(A))) & =f(g_\Lambda(A))=\iota_\Lambda(A).
    \end{align}
    Thus $f$ and $g$ are mutually inverse.
\end{proof}

\begin{lemma}[Evaluation of the algebraic equivalence]
    \label{lem:qca_compatible_algebraic_equiv_evaluation}
    \lean{SpinChain.CompatibleLocalAutomorphism.algebraicEquiv_apply,
        SpinChain.CompatibleLocalAutomorphism.algebraicEquiv_symm_apply}
    \leanok
    \uses{def:qca_compatible_algebraic_equiv, def:qca_compatible_local_homs}
    For every $A\in\mathcal A_{\mathrm{loc}}$,
    \begin{align}
        F(A) & =f(A), & F^{-1}(A) & =g(A).
    \end{align}
\end{lemma}

\begin{proof}\leanok
    \uses{def:qca_compatible_algebraic_equiv}
    Both identities follow from the defining forward and inverse maps of $F$.
\end{proof}

\begin{definition}[Forward local norm hypothesis]
    \label{def:qca_compatible_local_norm}
    \lean{SpinChain.CompatibleLocalAutomorphism.ForwardNormPreserving}
    \leanok
    \uses{def:qca_compatible_local_automorphism, def:qca_algebraic_local_operator_norm}
    Assume $d>0$.  The forward local norm hypothesis is
    \begin{align}
        \lVert f_\Lambda(A)\rVert & =\lVert A\rVert
    \end{align}
    for every finite region $\Lambda$ and every $A\in\mathcal A_\Lambda$.
\end{definition}

\begin{lemma}[Forward algebraic norm formula]
    \label{lem:qca_compatible_algebraic_norm}
    \lean{SpinChain.CompatibleLocalAutomorphism.norm_algebraicEquiv}
    \leanok
    \uses{def:qca_compatible_algebraic_equiv, def:qca_compatible_local_norm, def:qca_algebraic_local_operator_norm}
    Under the forward local norm hypothesis, every
    $A\in\mathcal A_{\mathrm{loc}}$ satisfies
    \begin{align}
        \lVert F(A)\rVert & =\lVert A\rVert.
    \end{align}
\end{lemma}

\begin{proof}\leanok
    \uses{lem:qca_compatible_algebraic_equiv_evaluation, lem:qca_compatible_local_homs_evaluation, def:qca_algebraic_local_operator_norm}
    By induction on finite-region representatives, for
    $A=\iota_\Lambda(A_\Lambda)$,
    \begin{align}
        \lVert F(A)\rVert
         & =\lVert f_\Lambda(A_\Lambda)\rVert
        =\lVert A_\Lambda\rVert
        =\lVert A\rVert.
    \end{align}
\end{proof}

\begin{lemma}[Forward and inverse algebraic isometries]
    \label{lem:qca_compatible_algebraic_isometry}
    \lean{SpinChain.CompatibleLocalAutomorphism.algebraicEquiv_isometry,
        SpinChain.CompatibleLocalAutomorphism.algebraicEquiv_symm_isometry}
    \leanok
    \uses{lem:qca_compatible_algebraic_norm, def:qca_compatible_algebraic_equiv}
    Under the forward local norm hypothesis, for all
    $A,B\in\mathcal A_{\mathrm{loc}}$,
    \begin{align}
        \lVert F(A)-F(B)\rVert           & =\lVert A-B\rVert, &
        \lVert F^{-1}(A)-F^{-1}(B)\rVert & =\lVert A-B\rVert.
    \end{align}
\end{lemma}

\begin{proof}\leanok
    \uses{lem:qca_compatible_algebraic_norm, def:qca_compatible_algebraic_equiv}
    Applying the forward norm identity to $A-B$ proves that $F$ is an isometry.
    Since $F\circ F^{-1}$ is the identity, the right-inverse law then makes
    $F^{-1}$ an isometry as well.
\end{proof}

\begin{lemma}[Inverse algebraic norm formula]
    \label{lem:qca_compatible_algebraic_inverse_norm}
    \lean{SpinChain.CompatibleLocalAutomorphism.norm_algebraicEquiv_symm}
    \leanok
    \uses{lem:qca_compatible_algebraic_isometry}
    Under the forward local norm hypothesis, every
    $A\in\mathcal A_{\mathrm{loc}}$ satisfies
    \begin{align}
        \lVert F^{-1}(A)\rVert & =\lVert A\rVert.
    \end{align}
\end{lemma}

\begin{proof}\leanok
    \uses{lem:qca_compatible_algebraic_isometry}
    Applying the inverse isometry to the distance from $A$ to $0$ and using
    $F^{-1}(0)=0$ gives the reverse norm identity.
\end{proof}

\begin{definition}[Completed compatible local automorphism]
    \label{def:qca_compatible_quasi_local_equiv}
    \lean{SpinChain.CompatibleLocalAutomorphism.quasiLocalEquiv}
    \leanok
    \uses{def:qca_compatible_algebraic_equiv, lem:qca_compatible_algebraic_isometry, def:completion_star_alg_equiv}
    Under the forward local norm hypothesis, the algebraic equivalence extends
    to a star-algebra equivalence
    \begin{align}
        \overline F\colon\mathcal A & \simeq\mathcal A
    \end{align}
    of the quasi-local completion.
\end{definition}

\begin{proof}\leanok
    \uses{lem:qca_compatible_algebraic_isometry, def:completion_star_alg_equiv}
    The forward norm formula makes $F$ an isometry, and the inverse law makes
    $F^{-1}$ an isometry, so the completion extension for star-algebra
    equivalences applies to these two continuous maps.
\end{proof}

\begin{lemma}[Dense evaluation of the completed automorphism]
    \label{lem:qca_compatible_quasi_local_dense_evaluation}
    \lean{SpinChain.CompatibleLocalAutomorphism.quasiLocalEquiv_algebraicToQuasiLocal,
        SpinChain.CompatibleLocalAutomorphism.quasiLocalEquiv_symm_algebraicToQuasiLocal}
    \leanok
    \uses{def:qca_compatible_quasi_local_equiv, def:completion_star_alg_equiv}
    Let $j\colon\mathcal A_{\mathrm{loc}}\to\mathcal A$ be the canonical
    completion map.  For every $A\in\mathcal A_{\mathrm{loc}}$,
    \begin{align}
        \overline F(j(A))      & =j(F(A)),      &
        \overline F^{-1}(j(A)) & =j(F^{-1}(A)).
    \end{align}
\end{lemma}

\begin{proof}\leanok
    \uses{def:qca_compatible_quasi_local_equiv, def:completion_star_alg_equiv}
    These are the canonical-image formulas for the continuous extension of a
    star-algebra equivalence and for its inverse.
\end{proof}

\begin{lemma}[Finite evaluation of the completed automorphism]
    \label{lem:qca_compatible_quasi_local_finite_evaluation}
    \lean{SpinChain.CompatibleLocalAutomorphism.quasiLocalEquiv_quasiLocalObservable,
        SpinChain.CompatibleLocalAutomorphism.quasiLocalEquiv_symm_quasiLocalObservable}
    \leanok
    \uses{lem:qca_compatible_quasi_local_dense_evaluation, lem:qca_compatible_algebraic_equiv_evaluation, lem:qca_compatible_local_homs_evaluation, def:qca_quasi_local_observable}
    For every finite region $\Lambda$ and $A\in\mathcal A_\Lambda$,
    \begin{align}
        \overline F(j_\Lambda(A))      & =j(f_\Lambda(A)), &
        \overline F^{-1}(j_\Lambda(A)) & =j(g_\Lambda(A)).
    \end{align}
\end{lemma}

\begin{proof}\leanok
    \uses{lem:qca_compatible_quasi_local_dense_evaluation, lem:qca_compatible_algebraic_equiv_evaluation, lem:qca_compatible_local_homs_evaluation}
    Substituting $j_\Lambda=j\circ\iota_\Lambda$ into the dense evaluation
    formulas and using
    $F(\iota_\Lambda(A))=f_\Lambda(A)$ and
    $F^{-1}(\iota_\Lambda(A))=g_\Lambda(A)$.
\end{proof}

\begin{definition}[Finite-local support bound]
    \label{def:qca_compatible_local_support_bound}
    \lean{SpinChain.CompatibleLocalAutomorphism.MapsSupportWithin}
    \leanok
    \uses{def:qca_compatible_local_automorphism, def:qca_supported_in, def:qca_region_sumset}
    For a finite neighborhood $\mathcal N\subset\mathbb Z$, the forward family
    has support bound $\mathcal N$ if
    \begin{align}
        f_\Lambda(A) & \in
        \iota_{\Lambda+\mathcal N}
        (\mathcal A_{\Lambda+\mathcal N})
    \end{align}
    for every finite region $\Lambda$ and every $A\in\mathcal A_\Lambda$.
\end{definition}

\begin{lemma}[Propagation from a finite-local support bound]
    \label{thm:qca_compatible_local_propagation}
    \lean{SpinChain.CompatibleLocalAutomorphism.propagatesWithin_quasiLocalEquiv}
    \leanok
    \uses{def:qca_compatible_local_support_bound, lem:qca_compatible_quasi_local_finite_evaluation, def:qca_propagates_within}
    If the forward family has support bound $\mathcal N$, then the completed
    automorphism $\overline F$ propagates within $\mathcal N$.
\end{lemma}

\begin{proof}\leanok
    \uses{def:qca_compatible_local_support_bound, lem:qca_compatible_quasi_local_finite_evaluation}
    For $X=j_\Lambda(A)$ supported in $\Lambda$, the local support bound gives
    $B\in\mathcal A_{\Lambda+\mathcal N}$ such that
    $f_\Lambda(A)=\iota_{\Lambda+\mathcal N}(B)$.  Then
    \begin{align}
        \overline F(X)
         & =j(f_\Lambda(A))
        =j_{\Lambda+\mathcal N}(B),
    \end{align}
    so the image is supported in $\Lambda+\mathcal N$.
\end{proof}

\begin{definition}[Finite-local unit-translation formula]
    \label{def:qca_compatible_local_unit_translation}
    \lean{SpinChain.CompatibleLocalAutomorphism.CommutesWithUnitTranslation}
    \leanok
    \uses{def:qca_compatible_local_automorphism, def:qca_algebraic_local_translation}
    The forward family commutes with unit translation if, for every finite
    region $\Lambda$ and every $A\in\mathcal A_\Lambda$,
    \begin{align}
        f_{\Lambda+1}(\tau_{1,\Lambda}(A))
         & =\tau_1(f_\Lambda(A)).
    \end{align}
\end{definition}

\begin{lemma}[Algebraic commutation with unit translation]
    \label{lem:qca_compatible_algebraic_unit_translation}
    \lean{SpinChain.CompatibleLocalAutomorphism.algebraicEquiv_commutes_unitTranslation}
    \leanok
    \uses{def:qca_compatible_local_unit_translation, def:qca_compatible_algebraic_equiv, def:qca_algebraic_local_translation}
    If the finite-local unit-translation formula holds, then every
    $A\in\mathcal A_{\mathrm{loc}}$ satisfies
    \begin{align}
        F(\tau_1(A)) & =\tau_1(F(A)).
    \end{align}
\end{lemma}

\begin{proof}\leanok
    \uses{def:qca_compatible_local_unit_translation, lem:qca_compatible_algebraic_equiv_evaluation, lem:qca_compatible_local_homs_evaluation, def:qca_algebraic_local_translation}
    By induction on finite-region representatives, for
    $A=\iota_\Lambda(A_\Lambda)$.  Then
    \begin{align}
        F(\tau_1(A))
         & =f_{\Lambda+1}(\tau_{1,\Lambda}(A_\Lambda))
        =\tau_1(f_\Lambda(A_\Lambda))
        =\tau_1(F(A)).
    \end{align}
\end{proof}

\begin{lemma}[Translation covariance from unit translation]
    \label{thm:qca_compatible_translation_covariant}
    \lean{SpinChain.CompatibleLocalAutomorphism.translationCovariant_quasiLocalEquiv}
    \leanok
    \uses{lem:qca_compatible_algebraic_unit_translation, lem:qca_compatible_quasi_local_dense_evaluation, def:qca_quasi_local_translation, thm:qca_translation_covariant_iff_commute_one}
    If the finite-local unit-translation formula holds, then the completed
    automorphism $\overline F$ is translation covariant.
\end{lemma}

\begin{proof}\leanok
    \uses{lem:qca_compatible_algebraic_unit_translation, lem:qca_compatible_quasi_local_dense_evaluation, def:qca_quasi_local_translation, thm:qca_translation_covariant_iff_commute_one}
    By Lemma~\ref{thm:qca_translation_covariant_iff_commute_one}, it suffices
    to prove that $\overline F$ commutes with translation by one site.  By continuity
    on the completion, on every $A\in\mathcal A_{\mathrm{loc}}$,
    \begin{align}
        \overline F(\overline\tau_1(j(A)))
         & =j(F(\tau_1(A)))
        =j(\tau_1(F(A)))
        =\overline\tau_1(\overline F(j(A))).
    \end{align}
    The equality set is closed because both composites are continuous, so the
    identity extends from the canonical dense subalgebra to the completion.
\end{proof}

\begin{lemma}[QCA induced by compatible finite-local actions]
    \label{thm:qca_compatible_local_automorphism}
    \lean{SpinChain.CompatibleLocalAutomorphism.isQCA_quasiLocalEquiv}
    \leanok
    \uses{def:qca_compatible_quasi_local_equiv, thm:qca_compatible_local_propagation, thm:qca_compatible_translation_covariant, def:qca}
    Let $d>0$.  Suppose the compatible forward family preserves the local
    operator norm, has a finite support bound, and satisfies the finite-local
    unit-translation formula.  Then the completed automorphism $\overline F$ is
    a one-dimensional QCA.
\end{lemma}

\begin{proof}\leanok
    \uses{thm:qca_compatible_local_propagation, thm:qca_compatible_translation_covariant}
    The support hypothesis gives a finite forward propagation neighborhood,
    while the unit-translation formula gives translation covariance.  These
    are exactly the two defining properties of a one-dimensional QCA.
\end{proof}

\subsection{Centers and matrix sizes of the support factors}
\label{sec:qca_support_factors}

Let $d>0$ and let $\omega$ be a nearest-neighbour star-automorphism of the
homogeneous quasi-local chain. No translation covariance is assumed.
The following is the center and factor argument in
\cite[lines~1276--1282]{Gross2012IndexTheory}, under the homogeneous scope
recorded in \cite{gap:gnvw12_site_dependent_adjacent_generation_scope}.

\begin{lemma}[Scalar centers of the support algebras]
    \label{lem:qca_support_factor_centers}
    \lean{SpinChain.PropagatesWithin.embeddedSupportAlgebra_isCentral,
        SpinChain.PropagatesWithin.evenSupportAlgebra_isCentral,
        SpinChain.PropagatesWithin.oddSupportAlgebra_isCentral}
    \notready
    \uses{thm:qca_site_indexed_support_algebras_commute,
        cor:qca_support_algebra_commutant,thm:qca_quasi_local_center_scalar}
    Each support algebra $\mathcal R_y$ has scalar center:
    $Z(\mathcal R_y)=\mathbb C I$.
\end{lemma}
\begin{proof}
    If $X\in Z(\mathcal R_y)$, then $X$ commutes with $\mathcal R_y$ by
    centrality and with every $\mathcal R_z$, $z\ne y$, by pairwise
    commutation. Norm-closed generation implies $[X,A]=0$ for every
    quasi-local observable $A$, so $X=cI$. Faithfulness of the finite-region
    inclusion gives the same scalar-center assertion in the finite matrix
    coordinates.
\end{proof}

\begin{corollary}[Positive matrix presentations]
    \label{cor:qca_support_factor_presentations}
    \lean{SpinChain.PropagatesWithin.exists_evenSupportAlgebra_matrix,
        SpinChain.PropagatesWithin.exists_oddSupportAlgebra_matrix,
        SpinChain.PropagatesWithin.exists_supportAlgebra_matrixSizes}
    \notready
    \uses{lem:qca_support_factor_centers}
    There is a family of positive integers $r(y)$ with star-algebra
    isomorphisms
    \begin{align}
        \mathcal R_{2x}   & \cong M_{r(2x)}(\mathbb C),   &
        \mathcal R_{2x+1} & \cong M_{r(2x+1)}(\mathbb C).
        \label{eq:qca_support_factor_presentations}
    \end{align}
    The integer $r(y)$ is the matrix size; the complex dimension of
    $\mathcal R_y$ is $r(y)^2$.
\end{corollary}
\begin{proof}
    Apply the full-matrix scalar-center consequence of the finite-dimensional
    star-algebra classification to each finite support algebra. This is the
    consequence of \cite[Proposition~Csform]{SchumacherWerner2004QCA}
    documented in \cite{gap:sw04_csform_matrix_scope}; no general direct-sum
    classification is needed here. Choose these presentations simultaneously
    in the even and odd families and combine their matrix sizes according to
    the parity of the site index.
\end{proof}

\begin{lemma}[Commutant of the even-pair images]
    \label{lem:qca_even_pair_image_commutant}
    \lean{SpinChain.commute_of_forall_commute_image_evenPair}
    \notready
    \uses{thm:qca_even_pair_algebraic_generation,def:qca_quasi_local_algebra}
    Let $\omega$ be any star-automorphism of the homogeneous quasi-local
    chain. If $Z$ commutes with $\omega(\mathcal A_{E_x})$ for every even
    pair $E_x$, then $Z$ commutes with every quasi-local observable.
    Commutation with the separate support factors is not assumed.
    This is the intermediate generation argument in the strict-inclusion
    step of \cite[lines~1276--1287]{Gross2012IndexTheory}, under the
    homogeneous scope of
    \cite{gap:gnvw12_site_dependent_adjacent_generation_scope}.
\end{lemma}
\begin{proof}
    The even-pair algebras generate the algebraic local algebra. Their
    images lie in the closed star-algebra of elements commuting with $Z$
    and $Z^*$. Density of the local union and continuity of $\omega$ extend
    this containment to its entire image. Surjectivity makes that image
    the full quasi-local algebra.
\end{proof}

\begin{lemma}[Relative commutant inside an adjacent support pair]
    \label{lem:qca_adjacent_support_relative_commutant}
    \lean{SpinChain.PropagatesWithin.quasiLocalObservable_evenPair_mem_adjacentSupportJoin,
        SpinChain.PropagatesWithin.adjacentSupportJoin_commute,
        SpinChain.PropagatesWithin.existsUnique_eq_smul_one_of_mem_adjacentSupportJoin_commute_image}
    \notready
    \uses{lem:qca_even_pair_image_commutant,
        thm:qca_site_indexed_support_algebras_commute,thm:qca_quasi_local_center_scalar}
    Put $\mathcal J_x=\mathcal R_{2x}\vee\mathcal R_{2x+1}$.
    Then $\omega(\mathcal A_{E_x})\subseteq\mathcal J_x$, the algebras
    $\mathcal J_x$ and $\mathcal J_y$ commute for $x\ne y$, and
    \begin{align}
        \{Z\in\mathcal J_x:[Z,\omega(\mathcal A_{E_x})]=0\}
         & =\mathbb C I.
        \label{eq:qca_adjacent_support_relative_commutant}
    \end{align}
    This is the relative-commutant argument preceding the first adjacent
    generation identity in \cite[lines~1276--1287]{Gross2012IndexTheory}.
\end{lemma}
\begin{proof}
    The two-sided support containment gives the first inclusion. The four
    support labels of distinct pairs are distinct, so their pairwise
    commutation extends to the two joins. If $Z\in\mathcal J_x$ commutes
    with its own image-pair algebra, it therefore commutes with every other
    image-pair algebra as well. Their commutant is the global center, which
    consists exactly of scalars.
\end{proof}


\subsection{The homogeneous positive-rational index}
\label{sec:qca_gnvw_index}

Let $d>0$ and let $\omega$ be a nearest-neighbour star-automorphism of the
homogeneous chain. The following uses the fixed nearest-neighbour cell
presentation. The site-dependent source setting is recorded separately in
\cite{gap:gnvw12_site_dependent_adjacent_generation_scope}.

\begin{definition}[The fixed-cut GNVW index]
    \label{def:qca_gnvw_index}
    \lean{SpinChain.PropagatesWithin.supportAlgebraMatrixSize,
        SpinChain.PropagatesWithin.supportAlgebraMatrixSize_spec,
        SpinChain.PropagatesWithin.supportAlgebraMatrixSize_neZero,
        SpinChain.PropagatesWithin.gnvwIndex}
    \leanok
    \uses{lem:qca_physical_pair_adjacent_generation}
    Choose the coherent positive matrix-size family $r(y)$ of the actual
    support factors. Define
    \begin{align}
        \operatorname{ind}(\omega) & =\frac{r(0)}{d}\in\mathbb Q_{>0}.
        \label{eq:qca_gnvw_index_fixed_cut}
    \end{align}
    Its codomain is the multiplicative group of positive rationals.
    This is the homogeneous specialization of equation~dimtransfer in
    \cite[lines~1317--1329]{Gross2012IndexTheory}.
\end{definition}

\begin{lemma}[All cut presentations of the same index]
    \label{thm:qca_gnvw_index_dimtransfer}
    \lean{SpinChain.PropagatesWithin.supportAlgebraMatrixSize_even_add_one,
        SpinChain.PropagatesWithin.supportAlgebraMatrixSize_even_eq_zero,
        SpinChain.PropagatesWithin.coe_gnvwIndex_eq_even,
        SpinChain.PropagatesWithin.coe_gnvwIndex_eq_odd,
        SpinChain.PropagatesWithin.coe_gnvwIndex_eq_next_even,
        SpinChain.PropagatesWithin.gnvwIndex_ratio_eq,
        SpinChain.PropagatesWithin.coe_gnvwIndex_eq_of_evenSupportAlgebra_equiv,
        SpinChain.PropagatesWithin.supportAlgebraMatrixSize_even_eq_of_equiv,
        SpinChain.PropagatesWithin.supportAlgebraMatrixSize_odd_eq_of_equiv}
    \leanok
    \uses{def:qca_gnvw_index}
    For every $x\in\mathbb Z$,
    \begin{align}
        \operatorname{ind}(\omega)
         & =\frac{r(2x)}{d}
        =\frac{d}{r(2x+1)}
        =\frac{r(2x+2)}{d}.
        \label{eq:qca_gnvw_index_dimtransfer}
    \end{align}
    These are presentations using the same ordered support family. No
    translation covariance is assumed. The value is independent of the
    integer cut and of the chosen full-matrix presentation at that cut.
    This is homogeneous equation~dimtransfer of
    \cite[lines~1317--1324]{Gross2012IndexTheory}.
\end{lemma}
\begin{proof}\leanok
    The two size products equal $d^2$. Cancelling the positive common odd
    size gives $r(2x)=r(2x+2)$. Integer induction in both directions makes
    the even size constant. Dividing the first product identity gives
    the reciprocal odd ratio. Matrix-algebra isomorphisms identify the
    sizes of any two presentations of the same actual factor.
\end{proof}

\begin{lemma}[Support factors of a full tensor factor]
    \label{lem:qca_kronecker_embedding_supports}
    \lean{Matrix.leftSupportAlgebra_range_leftKroneckerEmbed,
        Matrix.rightSupportAlgebra_range_leftKroneckerEmbed,
        Matrix.leftSupportAlgebra_range_rightKroneckerEmbed,
        Matrix.rightSupportAlgebra_range_rightKroneckerEmbed}
    \leanok
    \uses{def:qca_left_support_algebra,def:qca_right_support_algebra}
    Let $I$ and $J$ be finite index sets. Then
    \begin{align}
        \operatorname{Spp}_{\mathrm L}
        \bigl(M_I(\mathbb C)\otimes I_J\bigr)
                                       & =M_I(\mathbb C)
        \quad\text{if }J\neq\emptyset, &
        \operatorname{Spp}_{\mathrm R}
        \bigl(M_I(\mathbb C)\otimes I_J\bigr)
                                       & =\mathbb C I_J,\notag   \\
        \operatorname{Spp}_{\mathrm L}
        \bigl(I_I\otimes M_J(\mathbb C)\bigr)
                                       & =\mathbb C I_I,       &
        \operatorname{Spp}_{\mathrm R}
        \bigl(I_I\otimes M_J(\mathbb C)\bigr)
                                       & =M_J(\mathbb C)
        \quad\text{if }I\neq\emptyset.
        \label{eq:qca_kronecker_embedding_supports}
    \end{align}
    Only the full support of each factor needs the spectator index set to
    be nonempty. These are the support calculations in the shift example
    following equation~RR2x of
    \cite[lines~1251--1268]{Gross2012IndexTheory}.
\end{lemma}
\begin{proof}\leanok
    For $A\in M_I(\mathbb C)$, $B\in M_J(\mathbb C)$, $i,j\in I$ and
    $r,s\in J$, the coefficients are
    \begin{align}
        (A\otimes I_J)^{\mathrm L}_{rs}
         & =\delta_{rs}A,    &
        (A\otimes I_J)^{\mathrm R}_{ij}
         & =A_{ij}I_J,\notag   \\
        (I_I\otimes B)^{\mathrm L}_{rs}
         & =B_{rs}I_I,       &
        (I_I\otimes B)^{\mathrm R}_{ij}
         & =\delta_{ij}B.
        \label{eq:qca_kronecker_embedding_coefficients}
    \end{align}
    Taking $r=s$, which needs $J\neq\emptyset$, recovers every $A$ as a
    left coefficient; taking $i=j$, which needs $I\neq\emptyset$, recovers
    every $B$ as a right coefficient. The remaining coefficients are scalar
    multiples of the identity, so they generate only the scalars.
\end{proof}

\begin{lemma}[Finite-region configuration count]
    \label{lem:qca_finite_region_configuration_count}
    \lean{SpinChain.card_config}
    \leanok
    A region $\Lambda$ has $d^{|\Lambda|}$ configurations. In particular,
    a physical pair has Hilbert-space dimension $d^2$ and observable-algebra
    dimension $d^4$.
\end{lemma}
\begin{proof}\leanok
    There is one independent choice among $d$ values at each site.
\end{proof}

\begin{lemma}[Faithful bipartite observable coordinates]
    \label{lem:qca_bipartite_observable_embedding}
    \lean{SpinChain.bipartiteMatrixToQuasiLocalObservable,
        SpinChain.bipartiteMatrixToQuasiLocalObservable_injective,
        SpinChain.bipartiteMatrixToQuasiLocalObservable_leftKroneckerEmbed,
        SpinChain.bipartiteMatrixToQuasiLocalObservable_rightKroneckerEmbed,
        SpinChain.map_range_leftKroneckerEmbed_bipartiteMatrixToQuasiLocalObservable,
        SpinChain.map_range_rightKroneckerEmbed_bipartiteMatrixToQuasiLocalObservable}
    \leanok
    For disjoint finite regions, the canonical bipartite matrix coordinates
    embed faithfully into the quasi-local algebra. The full left and right
    factor ranges become exactly the observable algebras of their respective
    finite regions. These are the finite coordinates of
    \cite[lines~1251--1278]{Gross2012IndexTheory}.
\end{lemma}
\begin{proof}\leanok
    Write $\varphi$ for the bipartite embedding and
    $\iota_\Lambda,\iota_\Gamma$ for the canonical finite-region
    observable inclusions. The ordered bipartite coordinate equivalence
    gives, for $A\in M_{\Lambda}(\mathbb C)$ and
    $B\in M_{\Gamma}(\mathbb C)$,
    \begin{align}
        \varphi(A\otimes I_\Gamma)
         & =\iota_\Lambda(A), &
        \varphi(I_\Lambda\otimes B)
         & =\iota_\Gamma(B).
        \label{eq:qca_bipartite_observable_embedding}
    \end{align}
    Injectivity of $\varphi$ follows from faithfulness of the canonical
    inclusion of $\Lambda\cup\Gamma$. The two identities map the full left
    and right factor ranges onto the ranges of $\iota_\Lambda$ and
    $\iota_\Gamma$.
\end{proof}

\begin{lemma}[Direct shift normalization]
    \label{thm:qca_gnvw_index_shift}
    \lean{SpinChain.quasiLocalTranslation_propagatesWithin_singleton,
        SpinChain.quasiLocalTranslation_propagatesWithin_neg_one,
        SpinChain.quasiLocalTranslation_propagatesWithin_one,
        SpinChain.range_quasiLocalTranslation_comp_quasiLocalObservable,
        SpinChain.translateRegion_neg_one_evenPair,
        SpinChain.translateRegion_one_evenPair,
        SpinChain.PropagatesWithin.evenPairLocalImage_quasiLocalTranslation_neg_one,
        SpinChain.PropagatesWithin.evenPairLocalImage_quasiLocalTranslation_one,
        SpinChain.PropagatesWithin.supportAlgebras_quasiLocalTranslation_neg_one,
        SpinChain.PropagatesWithin.supportAlgebras_quasiLocalTranslation_one,
        SpinChain.PropagatesWithin.supportAlgebraMatrixSize_quasiLocalTranslation_neg_one_even,
        SpinChain.PropagatesWithin.supportAlgebraMatrixSize_quasiLocalTranslation_neg_one_odd,
        SpinChain.PropagatesWithin.supportAlgebraMatrixSize_quasiLocalTranslation_one_even,
        SpinChain.PropagatesWithin.supportAlgebraMatrixSize_quasiLocalTranslation_one_odd,
        SpinChain.PropagatesWithin.gnvwIndex_quasiLocalTranslation_neg_one,
        SpinChain.PropagatesWithin.gnvwIndex_quasiLocalTranslation_one}
    \leanok
    \uses{def:qca_gnvw_index,thm:qca_gnvw_index_dimtransfer,
        lem:qca_kronecker_embedding_supports,lem:qca_bipartite_observable_embedding,
        lem:qca_finite_region_configuration_count}
    Write $T_a$ for observable translation by $a$. The source Heisenberg
    shift is $T_{-1}$, which moves observables left and states right.
    Directly in the ordered finite coordinates,
    \begin{align}
        \mathcal R_{2x}(T_{-1})   & =\mathcal A_{L_x},  &
        \mathcal R_{2x+1}(T_{-1}) & =\mathbb C I,\notag   \\
        \mathcal R_{2x}(T_1)      & =\mathbb C I,       &
        \mathcal R_{2x+1}(T_1)    & =\mathcal A_{P_x}.
        \label{eq:qca_translation_support_factors}
    \end{align}
    Hence
    \begin{align}
        \operatorname{ind}(T_{-1}) & =d,      &
        \operatorname{ind}(T_1)    & =d^{-1}.
        \label{eq:qca_gnvw_index_shift}
    \end{align}
    The first is the homogeneous shift normalization of
    \cite[lines~1325--1329]{Gross2012IndexTheory}; the second is a direct
    orientation check. The repeated even subscript in the source scalar
    support example at line~1268 is corrected as documented in
    \cite{gap:gnvw12_shift_example_subscript}.
\end{lemma}
\begin{proof}\leanok
    Translation has its singleton propagation neighbourhood, which lies
    inside the nearest-neighbour interval. It maps the full source-pair
    algebra onto the left pair for $T_{-1}$, and onto the right pair for
    $T_1$. Faithful bipartite coordinates identify the finite images with
    the full left and right tensor factors respectively. Their supports
    therefore have the displayed full/scalar values. The full factor has
    matrix size $d^2$; the companion product identity makes the scalar
    factor's matrix size one. Substitution into the already derived index
    ratios proves both values without any composition or inversion law.
\end{proof}


\subsection{Unitary gates for the actual adjacent factors}
\label{sec:qca_margolus_unitary_gates}

Let $d>0$ and let $\omega$ be a nearest-neighbour star-automorphism of the
homogeneous chain. Write $E_x=\{2x,2x+1\}$,
$L_x=\{2x-1,2x\}$ and $P_x=\{2x+1,2x+2\}$.
Choose the already derived coherent matrix sizes $n(y)>0$ of the support
factors, with $n(2x)n(2x+1)=n(2x+1)n(2x+2)=d^2$.
The gates below may depend on $x$; no translation covariance is assumed.
These are the finite spatial steps of generalized Margolus partitioning
in \cite[equations~bruleU and bruleV, lines~980--1001]{SchumacherWerner2004QCA},
using the adjacent-generation equalities of
\cite[lines~1276--1318]{Gross2012IndexTheory}.

\begin{lemma}[Prescribed coordinates of commuting products]
    \label{lem:qca_commuting_product_coordinates}
    \lean{StarSubalgebra.commutingMulEquivSup,
        StarSubalgebra.commutingMulEquivSup_tmul_coe,
        StarSubalgebra.commutingSupEquivMatrix_apply_mul,
        SpinChain.PropagatesWithin.adjacentSupportJoinMatrixEquiv,
        SpinChain.PropagatesWithin.adjacentSupportJoinMatrixEquiv_apply_product,
        SpinChain.PropagatesWithin.physicalPairSupportJoinMatrixEquiv,
        SpinChain.PropagatesWithin.physicalPairSupportJoinMatrixEquiv_apply_product}
    \notready
    \uses{thm:qca_commuting_full_matrix_product,
        cor:qca_support_factor_presentations}
    For commuting full matrix factors $\mathcal R$ and $\mathcal S$ with
    prescribed presentations $e_R$ and $e_S$, multiplication identifies
    $\mathcal R\otimes\mathcal S$ with $\mathcal R\vee\mathcal S$ and
    the prescribed matrix coordinates satisfy
    \begin{align}
        e_{R\vee S}(ab) & =e_R(a)\otimes e_S(b).
    \end{align}
    This applies in order to the even and odd support factors, and to
    the odd and following even support factors.
\end{lemma}
\begin{proof}
    \notready
    Multiplication is injective by the full matrix factor theorem and
    has the join as its range. Transport the tensor product of the two
    presentations through this multiplication equivalence. A pure tensor
    maps to its actual ambient product and to its Kronecker coordinates.
\end{proof}

\begin{lemma}[The two actual matrix rules]
    \label{lem:qca_margolus_matrix_rules}
    \lean{SpinChain.PropagatesWithin.evenSupportMatrixEquiv,
        SpinChain.PropagatesWithin.oddSupportMatrixEquiv,
        SpinChain.PropagatesWithin.margolusInputRule,
        SpinChain.PropagatesWithin.margolusOutputRule,
        SpinChain.PropagatesWithin.margolusInputRule_symm_kronecker,
        SpinChain.PropagatesWithin.margolusOutputRule_kronecker}
    \notready
    \uses{lem:qca_commuting_product_coordinates,
        thm:qca_image_pair_adjacent_generation,
        thm:qca_physical_pair_adjacent_generation}
    The actual even-pair image and the actual shifted physical-pair join
    give star-isomorphisms
    \begin{align}
        F_x & : M_{d^2}(\mathbb C)\longrightarrow
        M_{n(2x)n(2x+1)}(\mathbb C),\notag                   \\
        G_x & : M_{n(2x+1)n(2x+2)}(\mathbb C)\longrightarrow
        M_{d^2}(\mathbb C).
    \end{align}
    If $\iota_\Lambda$ denotes canonical finite-region inclusion, then
    for $a\in\mathcal R_{2x}$ and $b\in\mathcal R_{2x+1}$,
    \begin{align}
        \omega\bigl(\iota_{E_x}(F_x^{-1}(e_{2x}(a)\otimes e_{2x+1}(b)))\bigr)
         & =\iota_{L_x}(a)\iota_{P_x}(b).
    \end{align}
    For $b\in\mathcal R_{2x+1}$ and $c\in\mathcal R_{2x+2}$,
    \begin{align}
        \iota_{P_x}(G_x(e_{2x+1}(b)\otimes e_{2x+2}(c)))
         & =\iota_{P_x}(b)\iota_{L_{x+1}}(c).
    \end{align}
\end{lemma}
\begin{proof}
    \notready
    Restrict the injective even-pair image homomorphism to its actual
    range, which is its two-factor join. Compose with its prescribed
    matrix presentation. For the shifted rule, use the inverse of the
    shifted join presentation and the inverse of faithful physical-pair
    inclusion. The preceding product-coordinate identity gives both
    displayed formulas.
\end{proof}

\begin{lemma}[Spatial unitaries for both local rules]
    \label{lem:qca_margolus_spatial_unitaries}
    \lean{SpinChain.PropagatesWithin.exists_margolus_unitary_gates}
    \notready
    \uses{lem:qca_margolus_matrix_rules,lem:qca_full_matrix_spatial_implementer}
    At every integer position $x$ there are unitary maps
    \begin{align}
        U_x & : \mathbb C^{d^2}\longrightarrow
        \mathbb C^{n(2x)}\otimes\mathbb C^{n(2x+1)},\notag    \\
        V_x & : \mathbb C^{n(2x+1)}\otimes\mathbb C^{n(2x+2)}
        \longrightarrow\mathbb C^{d^2}
    \end{align}
    such that $F_x(A)=U_xAU_x^*$ and $G_x(B)=V_xBV_x^*$.
    Consequently the two preceding local identities hold with
    $F_x^{-1}(B)=U_x^*BU_x$ and $G_x(B)=V_xBV_x^*$.
\end{lemma}
\begin{proof}
    \notready
    Apply finite full matrix spatial implementation to each actual
    matrix rule. The unitary identities give the inverse conjugation
    formula. Substitute the prescribed factor products into the two
    actual local-rule identities.
\end{proof}

\begin{lemma}[Cancelling intermediate basis freedom]
    \label{lem:qca_margolus_basis_freedom}
    \lean{SpinChain.PropagatesWithin.margolus_basisChange_unitary,
        SpinChain.PropagatesWithin.margolus_basisChange_input,
        SpinChain.PropagatesWithin.margolus_basisChange_output}
    \notready
    \uses{lem:qca_margolus_spatial_unitaries}
    Let $P,Q$ be independent unitary basis changes on two intermediate
    factors, and write $K=P\otimes Q$. The changed input and output gates
    $U'=KU$ and $V'=VK^*$ are unitary whenever $U,V$ are unitary. For all
    factor observables $A,B$,
    \begin{align}
        U'^*((PAP^*)\otimes(QBQ^*))U' & =U^*(A\otimes B)U,\notag \\
        V'((PAP^*)\otimes(QBQ^*))V'^* & =V(A\otimes B)V^*.
    \end{align}
    Thus a basis rotation on each actual intermediate factor cancels
    between its input and output appearances. This is the freedom in
    \cite[lines~1033--1041]{SchumacherWerner2004QCA}; it imposes no invented
    translation constraint on an arbitrary pair of gates.
\end{lemma}
\begin{proof}
    \notready
    Kronecker products preserve both unitary identities. Substitution of
    the changed gates gives $K^*K=I$ on both sides of the intermediate
    observable; cancellation proves the two identities.
\end{proof}


\subsection{The actual QCA induced by an MPU}
\label{sec:mpu_actual_qca}

\begin{lemma}[Physical-adjoint finite conjugation is the actual inverse]
    \label{lem:mpu_actual_finite_adjoint_inverse}
    \lean{SpinChain.finiteChainMPOMatrix_physicalAdjointTensor,
        SpinChain.finiteChainConjugation_physicalAdjointTensor_apply_apply,
        SpinChain.finiteChainConjugation_apply_physicalAdjointTensor_apply,
        SpinChain.finiteChainConjugation_physicalAdjointTensor}
    \notready
    \uses{def:mpu_finite_chain_conjugation}
    For every MPU and $N>1$, physical adjunction gives the adjoint
    periodic operator on the same anchored finite chain. Its actual
    finite conjugation is the inverse of the original finite
    conjugation. This is the unitary inverse identity underlying
    \cite[lines~2300--2306]{Cirac2017MPU}; no inverse-local-family
    condition is assumed.
\end{lemma}
\begin{proof}
    Physical adjunction interchanges the physical matrix indices and
    conjugates coefficients, so it gives $U^{(N)\dagger}$ without
    reversing site order. The two actual finite conjugations cancel
    by the two unitary equations.
\end{proof}

\begin{lemma}[Actual QCA from canonical form II]
    \label{lem:mpu_actual_qca_canonical}
    \lean{SpinChain.exists_isQCA_of_isMPUCanonicalFormII}
    \notready
    \uses{lem:mpu_actual_finite_adjoint_inverse,
        thm:mpu_finite_heisenberg_completed_forward_action,
        lem:mpu_finite_heisenberg_translation_covariance,
        lem:mpu_finite_heisenberg_forward_propagation}
    A canonical-form-II MPU of bond dimension $D$ determines an actual
    quasi-local QCA with propagation neighbourhood
    $[-4D^4,4D^4]\cap\mathbb Z$. Its action is the actual original-site
    finite MPU conjugation on every interval of length greater than
    one containing the input support and its complete neighbourhood.
    The concrete original and physical-adjoint standard forms are
    derived from the canonical datum. This is the canonical
    intermediate construction in
    \cite[lines~378--427, 563--622 and~2292--2306]{Cirac2017MPU}.
\end{lemma}
\begin{proof}
    The established simple-block construction supplies a positive block
    length $k\le D^2$ and its actual standard form. Physical adjunction
    preserves the canonical datum, and the same construction supplies
    an adjoint block and standard form independently. Their completed
    forward maps have derived finite support. For a finite input,
    choose one common chain buffering both its support and the actual
    first output support. Faithfulness of the finite-region inclusion
    identifies that output with the actual first finite conjugation.
    The two finite conjugations cancel by the preceding lemma.
    Density extends both inverse identities to the completion. The
    resulting star-automorphism has the derived covariance and
    propagation properties. Finally $D^2\le D^4$ for $D>0$ gives the
    stated neighbourhood, and the actual finite-chain evaluation
    holds on every fully buffered interval.
\end{proof}

\begin{theorem}[The actual QCA of a matrix product unitary]
    \label{thm:mpu_actual_qca}
    \lean{SpinChain.exists_isQCA_of_isMPU,
        SpinChain.mpuQuasiLocalEquiv,
        SpinChain.isQCA_mpuQuasiLocalEquiv,
        SpinChain.mpuQuasiLocalEquiv_propagatesWithin,
        SpinChain.mpuQuasiLocalEquiv_finiteChainConjugation,
        SpinChain.mpuQuasiLocalEquiv_apply_eq_bufferedMPUQuasiLocalAction}
    \notready
    \uses{lem:mpu_actual_qca_canonical}
    Let $d>0$ and let $U$ be an MPU of bond dimension $D$.
    There is an actual quasi-local star-automorphism $\Phi_U$ which is
    translation covariant and propagates within
    \[
        \mathcal N_D=[-4D^4,4D^4]\cap\mathbb Z.
    \]
    For every finite support $\Lambda$, every anchored interval
    $I_{b,N}=[b,b+N)\cap\mathbb Z$ with $N>1$, and every
    $A\in\mathcal A_\Lambda$, if $\Lambda$ and
    $\Lambda+\mathcal N_D$ lie in $I_{b,N}$, then
    \[
        \Phi_U(A)=U_b^{(N)}\,\iota_{\Lambda,I_{b,N}}(A)\,
        U_b^{(N)\dagger}
    \]
    as a quasi-local observable. Thus the finite-chain expression is
    independent of the sufficiently buffered anchor and length,
    with no divisibility condition. The same action is obtained from
    any concrete standard form of a positive block, by comparing the
    actual finite conjugations on a common buffer. The only data are positive physical
    dimension and the actual MPU property; bond positivity is derived.
    This is the original-site construction and finite-chain formula of
    \cite[lines~2292--2306]{Cirac2017MPU}. The historical scope record
    and pending verification are in
    \cite{gap:mpu_qca_complete_block_heisenberg_scope}.
\end{theorem}
\begin{proof}
    The existing reduced canonical representative has bond dimension
    $0<D_{\mathrm{red}}\le D$ and exactly the same periodic physical
    operator at every positive length. Apply the preceding construction
    to it. Enlarge its neighbourhood using
    $4D_{\mathrm{red}}^4\le4D^4$. For every $N>1$, equality of the
    actual periodic physical matrices identifies the two anchored
    finite conjugations, so the resulting QCA has the original
    tensor's finite-chain formula. No source factors, selected cut
    spaces or virtual gauges are transported between the possibly
    different bond dimensions.
\end{proof}

